\documentclass[10pt,a4paper]{article}

% ----------------- Paquetes -------------------%
\usepackage[T1]{fontenc}
\usepackage[utf8]{inputenc}
\usepackage[a4paper,margin=2.5cm]{geometry}
\usepackage{mathptmx}             
\usepackage{changepage} 
\usepackage{setspace}
\usepackage[spanish]{babel}
\usepackage[labelfont=bf,labelsep=space]{caption}
\usepackage{amssymb}
\usepackage{amsmath}
\usepackage{ebgaramond}
\usepackage{placeins}
\usepackage{etoolbox}
\usepackage{setspace}
\usepackage{parskip} 
\usepackage{tcolorbox}
\usepackage{needspace}
\usepackage{xparse}
\usepackage{ocgx2}
\usepackage{hyperref}
\usepackage{hypcap}
\usepackage{soul}\sethlcolor{yellow}% resaltado amarillo: \hl{texto} (Panel TCEE, ⌃H)



%%%%%%%%%%%%%%%%%%%%%%%%%%%%%%%%%%%%%%%%%%%%%%%%%%%%%%%%%%%%%%%%
% --- Fija primero tus valores globales "buenos" ---
\setstretch{1.2}                 % Interlineado global
\setlength{\parskip}{0.7\baselineskip}
\setlength{\parindent}{0pt}
\newlength{\GlobalParskip}
\setlength{\GlobalParskip}{\parskip}

\newcounter{eqnum}[subsection]
\renewcommand{\theeqnum}{\arabic{eqnum}}
\newcounter{alphaitem}
\newcounter{numitem}

%% ------------------- Recuadros----------------------%%
\newcommand{\puntosclave}[1]{%
  \begin{tcolorbox}[
    colframe=black,
    title={Puntos clave},
    before upper={\setlength{\parindent}{0.5cm}\setlength{\parskip}{0.5\baselineskip}}
  ]
  #1
  \end{tcolorbox}
}

\newcommand{\modificaciones}[1]{
  \begin{tcolorbox}[
    colback=white,
    colframe=red,
    title={Modificaciones a realizar},
    before upper={\setlength{\parindent}{0.5cm}\setlength{\parskip}{0.5\baselineskip}}
  ]
  #1
  \end{tcolorbox}
}

%% ------------------- Preguntas TEST----------------------%%
% Opciones: radio (single-choice)
\NewDocumentCommand{\OptRadio}{m m m}{%
  \noindent\CheckBox[radio,name=#1,value=#2,width=1.2ex,height=1.2ex]{}%
  \;\textbf{#2.}~#3\par
}

% Opciones: checkbox (multi-choice)
\NewDocumentCommand{\OptCheck}{m m}{%
  \noindent\CheckBox[name=#1,width=1.2ex,height=1.2ex,bordercolor={0 0 0}]{}%
  \;#2\par
}

% Pregunta single-choice (radio)
% Args: 1=id, 2=enunciado, 3=opciones, 4=solución+justificación, 5=imagen (o {}), 6=origen
\NewDocumentCommand{\PreguntaTestSingle}{m m m m m m}{%
  \par\medskip
  \noindent\textbf{#1.}~#2\par\smallskip

    % Imagen (si no es {})
    \IfBlankTF{#5}{}{%
      \begin{center}
        \includegraphics[width=0.75\linewidth]{#5}
      \end{center}
    }

    \begin{Form}
      #3
    \end{Form}

    \par\smallskip
    \switchocg{sol-#1}{\fbox{Mostrar / ocultar solución}}%
    \hfill{\footnotesize #6}\par\smallskip

    \begin{ocg}{Solución}{sol-#1}{0}
      \noindent #4
    \end{ocg}
  \par\medskip
}

% Pregunta multi-choice (checkboxes)
\NewDocumentCommand{\PreguntaTestMulti}{m m m m m m}{%
  \par\medskip
  \noindent\textbf{#1.}~#2\par\smallskip

    % Imagen (si no es {})
    \IfBlankTF{#5}{}{%
      \begin{center}
        \includegraphics[width=0.75\linewidth]{#5}
      \end{center}
    }
    \begin{Form}
      #3
    \end{Form}

    \par\smallskip
    \switchocg{sol-#1}{\fbox{Mostrar / ocultar solución}}%
    \hfill{\footnotesize #6}\par\smallskip

    \begin{ocg}{Solución}{sol-#1}{0}
      \noindent #4
    \end{ocg}
  \par\medskip
}

%% ------------------- Listas----------------------%%
    % --- Lista alfabética ---
    \newenvironment{la}
      {\begin{list}{\alph{alphaitem})}{%
          \usecounter{alphaitem}%
          \setlength{\leftmargin}{1cm}%
          \setlength{\parskip}{3pt}% 
          \setlength{\itemsep}{0pt}% ← espacio entre ítems
          \setlength{\parsep}{0pt}%  ← párrafos dentro de un ítem
      }}
      {\end{list}}
      
    % --- Lista numérica ---
    \newenvironment{lnum}
      {\begin{list}{\arabic{numitem}.}{%
          \usecounter{numitem}%
          \setlength{\leftmargin}{1cm}%
          \setlength{\parskip}{3pt}% 
          \setlength{\itemsep}{0pt}% ← espacio entre ítems
          \setlength{\parsep}{0pt}%  ← párrafos dentro de un ítem
      }}
      {\end{list}}

%% ------------------- Preguntas Test ----------------------%%
      \usepackage{xparse} % si ya lo tienes, no lo repitas

    \NewDocumentEnvironment{test}{m m o}{%
      \begin{tcolorbox}[
        colback=white,
        colframe=black!15,
        boxrule=0.6pt,
        arc=2mm,
        left=2mm,right=2mm,top=2mm,bottom=2mm
      ]
      \centering
      \includegraphics[width=\linewidth]{#1}
      \end{tcolorbox}
    
      \vspace{2mm}
    
      % Caja SOLO para texto (SÍ partible y mantiene márgenes al partir)
      \begin{tcolorbox}[
        enhanced,
        breakable,
        colback=black!3,
        colframe=black!25,
        boxrule=0.6pt,
        arc=2mm,
        left=2mm,right=2mm,top=1.5mm,bottom=1.5mm
      ]
      \textbf{Respuesta:} \textbf{#2}%
      \IfValueT{#3}{\par\smallskip \textbf{Justificación:} #3}
      \end{tcolorbox}
    }{}
      
% ------------------- Imágenes----------------------%
    \graphicspath{{Img/}}% Rutas por defecto 
    % --- Imagen adaptativa ---
    % Registros de dimensión definidos una sola vez
    \newdimen\imgcap
    \newdimen\imgmin
    \newdimen\imgmax
    \newdimen\imgremain
    \newdimen\imgh
    \newdimen\imgneed
    
    \newcommand{\imagenfit}[3]{%
      \addvspace{10pt}% separación antes
      \par\begingroup
        % Parámetros de composición (asignaciones, no \newdimen)
        \imgcap=1\baselineskip            % alto estimado de título+fuente
        \imgmin=0.15\textheight
        \imgmax=0.15\textheight
        \imgremain=\pagegoal
        \ifdim\pagegoal=\maxdimen \imgremain=0pt\fi
        \advance\imgremain by -\pagetotal
        \advance\imgremain by -\imgcap
        % Decisión de altura destino
        \ifdim\imgremain<\imgmin
          \newpage
          \imgh=0.20\textheight
        \else
          \imgh=\imgremain
          \ifdim\imgh>\imgmax \imgh=\imgmax \fi
        \fi
        % Reservar espacio para evitar cortes del bloque completo
        \imgneed=\imgh \advance\imgneed by \imgcap
        \Needspace{\imgneed}
        % Cuerpo
        \noindent\begin{minipage}{\textwidth}
          \centering
          \captionof{figure}{\textemdash\ #2}\par\vspace{0.3em}% título arriba
          \includegraphics[width=\linewidth,height=\imgh,keepaspectratio]{\detokenize{#1}}\par
          \vspace{0.3em}\small\textit{Fuente: }#3% fuente abajo
        \end{minipage}%
      \endgroup
      \par\addvspace{10pt}%
    }

%------------ Bloque de RECUADRO ESPECIAL -------------%
    \newenvironment{specialblock}[1]{%
      \begin{quote} % crea sangría a izquierda y derecha
      \small % texto más pequeño
      \noindent\textbf{#1}\\[-0.3em] % título en negrita
      \rule{\linewidth}{0.4pt}\\[0.6em]}{
      \end{quote}}
      
%-------------------------- Portada -----------------------%
\newcommand{\oldtitle}[1]{\def\OldTitle{#1}}
\newcommand{\changerationale}[1]{\def\ChangeWhy{#1}}

\newcommand{\changebox}{%
\begin{tcolorbox}[title={Historial de cambios del tema}]
\textbf{Título anterior:} \OldTitle\par
\textbf{Motivo del cambio:} \ChangeWhy
\end{tcolorbox}
}

%-------------------------- Author Font -----------------------%
    \newcommand{\authorfont}[1]{{\fontfamily{EBGaramond-TLF}\selectfont #1}}

%-------------------------- Anexos -----------------------%
    \makeatletter
    \newcounter{anexo}
    \renewcommand{\theanexo}{\Roman{anexo}}
    \newcommand{\anexo}[1]{%
      \clearpage
      \phantomsection
      \refstepcounter{anexo}%
      \noindent\textbf{Anexo \theanexo.- }#1\par\medskip
      \addcontentsline{stc}{section}{Anexo \theanexo.- #1}%
    }
    \makeatother

    %-------------------------- Lista de figuras (personal) -----------------------%
\makeatletter
\newcommand{\MyListOfFigures}{%
  \clearpage
  \phantomsection
  {\centering\bfseries\Large Índice de figuras\par}%
  \medskip
  % Entrada en nuestro índice de contenido (stc)
  \addcontentsline{stc}{section}{Índice de figuras}%
  % Recuperación del listado (sin cabecera propia)
  \begingroup
    \parindent=0pt\parskip=2pt
    \@starttoc{lof}%
  \endgroup
}
\makeatother

%-------------------------- Bibliografía -----------------------%
\makeatletter
% Contador para las referencias
\newcounter{myref}
% Cabecera de la bibliografía, entra en el índice 'stc'
\newcommand{\MyBibliography}{%
  \clearpage
  \phantomsection
  {\centering\bfseries\Large Bibliografía y referencias\par}%
  \medskip
  \addcontentsline{stc}{section}{Bibliografía y referencias}%
  \setcounter{myref}{0}%
}
\newcommand{\MyBibItem}[4]{%
  \refstepcounter{myref}%
  \noindent[\themyref]\ #1.\ \emph{#2}. #3. #4\par
}
\makeatother

%--------- Establecer cuatro niveles de índice --------%
    \setcounter{tocdepth}{4}   
    \setcounter{secnumdepth}{4}% numera hasta \paragraph

%%%%%%%%%%%%%%%%%%%%%%%%%%%%%%%%

\makeatletter

    %--------- Interlineado Global --------%
    \edef\GlobalBaseLineStretch{\baselinestretch}
    
    %--------- Espaciado de seccinones --------%
    \renewcommand\section{\@startsection{section}
      {1}{\z@}
      {2.5ex \@plus 1ex \@minus .2ex} % espacio antes
      {1.5ex \@plus .2ex}             % espacio después
      {\normalfont\Large\bfseries}    % formato del título
    }
    \renewcommand\subsection{\@startsection{subsection}{2}{\z@}
      {3ex \@plus 0.8ex \@minus .2ex}
      {1.5ex \@plus .2ex}
      {\normalfont\large\bfseries}
    }
    \renewcommand\subsubsection{\@startsection{subsubsection}{3}{\z@}
      {3ex \@plus 0.5ex \@minus .2ex}
      {1ex \@plus .2ex}
      {\normalfont\normalsize\bfseries}
    }
    \renewcommand\paragraph{\@startsection{paragraph}{4}{\z@}%
    {-2ex \@plus -0.5ex \@minus -.2ex}% espacio antes
    {0.8ex \@plus .2ex}% espacio después
    {\normalfont\normalsize\itshape}}% ← cursiva en lugar de negrita

    %--------- Ecuaciones --------%   
    % Variables globales para almacenar títulos y notas
    \gdef\leq@current@section{}%
    \gdef\leq@current@subsection{}%
    \gdef\leq@last@printed@section{}%
    \gdef\leq@last@printed@subsection{}%
    
    % Comando para añadir nota (invisible en el cuerpo)
    \newcommand{\eqnote}[1]{%
      \addcontentsline{leq}{leqnote}{#1}%
    }
    
    % Comando para ecuaciones
    \newcommand{\eqblock}[2]{%
      % Verificar si necesitamos imprimir el título de sección
      \ifx\leq@current@section\leq@last@printed@section\else
        \addcontentsline{leq}{leqsection}{\leq@current@section}%
        \global\let\leq@last@printed@section\leq@current@section
        \gdef\leq@last@printed@subsection{}% Reset subsección al cambiar sección
      \fi
      % Verificar si necesitamos imprimir el título de subsección
      \ifx\leq@current@subsection\leq@last@printed@subsection\else
        \ifx\leq@current@subsection\@empty\else
          \addcontentsline{leq}{leqsubsection}{\leq@current@subsection}%
          \global\let\leq@last@printed@subsection\leq@current@subsection
        \fi
      \fi
      % Ecuación
      \refstepcounter{eqnum}%
      \begin{equation*}
        #1\tag{E.\theeqnum}
      \end{equation*}%
      \addcontentsline{leq}{leqt}{[E.\theeqnum] -- #2}%
      \addcontentsline{leq}{leqm}{#1}%
    }
    
    %%% ----- Índice de ecuaciones ----------%%%
    % Tamaño para []
    \newcommand{\leqIndexFont}{\normalsize}
    \newcommand{\leqTitleFont}{\tiny}
    \newcommand{\leqNoteFont}{\tiny}
    % Cabecera de sección 
    \newcommand*\l@leqsection[2]{%
      \addvspace{25pt}% Mayor separación antes de nueva sección
      \noindent\leqIndexFont
      \makebox[\linewidth]{\rule{.38\linewidth}{0.4pt}\ \ \textbf{  \MakeUppercase{#1}}\ \ \rule{.38\linewidth}{0.4pt}}%
      \par\vspace{-10pt}
    }
    % Cabecera de subsección 
    \newcommand*\l@leqsubsection[2]{%
      \par\addvspace{15pt}%
      \noindent\leqIndexFont #1
      \par\vspace{-7pt}
    }
    % Título de ecuación 
    \newcommand*\l@leqt[2]{%
    \par\addvspace{10pt}%
      \par\noindent\leqTitleFont #1\par
    }
    % Miniatura de ecuación
    \newcommand*\l@leqm[2]{%
      \vspace{0.3\baselineskip}%
      \par\scriptsize\noindent\hspace*{1.5em}\(#1\)\par%
    }
    % Nota de ecuación 
    \newcommand*\l@leqnote[2]{%
      \vspace{0.2\baselineskip}%
      \par\noindent\leqNoteFont\hspace*{1.5em}\textbf{Nota.--} #1\par\vspace{0.5\baselineskip}%
    }
    % Generación del índice
    \newcommand{\mylistofequations}{%
      \clearpage
      \phantomsection
      % Título visual de la lista (ajústalo a tu gusto):
      {\centering\bfseries\Large Índice de ecuaciones\par}
      % Entrada en nuestro índice 'stc':
      \addcontentsline{stc}{section}{Índice de ecuaciones}%
      % Llamada a tu lista real (usa el nombre exacto de tu macro):
      \@starttoc{leq}%
    }
    
    %--------- Captura de títulos de sección --------%
    \let\leq@original@section\section
    \let\leq@original@subsection\subsection
    % Flag para saber si estamos en una sección asterisco
    \newif\ifLeqStarSection
    \LeqStarSectionfalse
    
    % Redefinir \section CON CONTROL DE ASTERISCO
    \renewcommand{\section}{%
      \@ifstar{\leq@section@star}{\@dblarg\leq@section@nostar}%
    }
    \def\leq@section@star#1{%
      \global\LeqStarSectiontrue
      \gdef\leq@current@section{#1}%
      \gdef\leq@current@subsection{}%
      \leq@original@section*{#1}%
      \global\LeqStarSectionfalse
    }
    \def\leq@section@nostar[#1]#2{%
      \global\LeqStarSectionfalse
      \gdef\leq@current@section{#2}%
      \gdef\leq@current@subsection{}%
      \leq@original@section[#1]{#2}%
    }
    % Redefinir \subsection
    \renewcommand{\subsection}{%
      \@ifstar{\leq@subsection@star}{\@dblarg\leq@subsection@nostar}%
    }
    \def\leq@subsection@star#1{%
      \global\LeqStarSectiontrue
      \gdef\leq@current@subsection{#1}%
      \leq@original@subsection*{#1}%
      \global\LeqStarSectionfalse
    }
    \def\leq@subsection@nostar[#1]#2{%
      \global\LeqStarSectionfalse
      \gdef\leq@current@subsection{#2}%
      \leq@original@subsection[#1]{#2}%
    }

    % ----- Restauración de valores Globales de estilo ----------%
    % Flags
    \newif\ifInFootnote
    \InFootnotefalse
    \pretocmd{\@footnotetext}{\InFootnotetrue}{}{}
    \apptocmd{\@footnotetext}{\InFootnotefalse}{}{}
    % Profundidad de listas (soporta anidadas)
    \newcount\MyListDepth \MyListDepth=0
    \AtBeginEnvironment{la}{\advance\MyListDepth by 1}
    \AfterEndEnvironment{la}{\advance\MyListDepth by -1}
    \AtBeginEnvironment{lnum}{\advance\MyListDepth by 1}
    \AfterEndEnvironment{lnum}{\advance\MyListDepth by -1}
    \AtBeginEnvironment{itemize}{\advance\MyListDepth by 1}
    \AfterEndEnvironment{itemize}{\advance\MyListDepth by -1}
    \AtBeginEnvironment{enumerate}{\advance\MyListDepth by 1}
    \AfterEndEnvironment{enumerate}{\advance\MyListDepth by -1}
    % Restauración diferida: hazlo al INICIO del próximo párrafo real
    \newif\if@pendingRestore
    \@pendingRestorefalse
    \newcommand{\RequestRestore}{\global\@pendingRestoretrue}
    \everypar{%
      \if@pendingRestore
        \global\@pendingRestorefalse
        \ifInFootnote\else
          \ifnum\MyListDepth=0
            \setlength{\parskip}{\GlobalParskip}%
            \setstretch{\GlobalBaseLineStretch}%
          \fi
        \fi
      \fi
    }
    % Al final de bloques:
    \AfterEndEnvironment{equation}{\RequestRestore}
    \AfterEndEnvironment{equation*}{\RequestRestore}
    \AfterEndEnvironment{la}{\RequestRestore}
    \AfterEndEnvironment{lnum}{\RequestRestore}
    % Al final de macros:
    \apptocmd{\eqblock}{\RequestRestore}{}{}
    \apptocmd{\imagenfit}{\RequestRestore}{}{}
    
\makeatother

% ===================== ToC propio con 4 niveles (único bloque) =====================
\makeatletter

% --- Escritores: añadimos una línea al .stc en cada cabecera numerada ---
% Guardamos las versiones actuales para llamar al comportamiento normal
\let\MyTOC@orig@section\section
\let\MyTOC@orig@subsection\subsection
\let\MyTOC@orig@subsubsection\subsubsection
\let\MyTOC@orig@paragraph\paragraph

% SECTION
\renewcommand{\section}{%
  \@ifstar{\MyTOC@section@star}{\@dblarg\MyTOC@section@nostar}%
}
\def\MyTOC@section@star#1{%
  \MyTOC@orig@section*{#1}% sin entrada al .stc
}
\def\MyTOC@section@nostar[#1]#2{%
  \MyTOC@orig@section[{#1}]{#2}% comportamiento normal (escribe al .toc)
  % Escribimos también nuestra línea al .stc (texto con \numberline)
  \addcontentsline{stc}{section}{\protect\numberline{\thesection}#2}%
}

% SUBSECTION
\renewcommand{\subsection}{%
  \@ifstar{\MyTOC@subsection@star}{\@dblarg\MyTOC@subsection@nostar}%
}
\def\MyTOC@subsection@star#1{%
  \MyTOC@orig@subsection*{#1}%
}
\def\MyTOC@subsection@nostar[#1]#2{%
  \MyTOC@orig@subsection[{#1}]{#2}%
  \addcontentsline{stc}{subsection}{\protect\numberline{\thesubsection}#2}%
}

% SUBSUBSECTION
\renewcommand{\subsubsection}{%
  \@ifstar{\MyTOC@subsubsection@star}{\@dblarg\MyTOC@subsubsection@nostar}%
}
\def\MyTOC@subsubsection@star#1{%
  \MyTOC@orig@subsubsection*{#1}%
}
\def\MyTOC@subsubsection@nostar[#1]#2{%
  \MyTOC@orig@subsubsection[{#1}]{#2}%
  \addcontentsline{stc}{subsubsection}{\protect\numberline{\thesubsubsection}#2}%
}

% PARAGRAPH
\renewcommand{\paragraph}{%
  \@ifstar{\MyTOC@paragraph@star}{\@dblarg\MyTOC@paragraph@nostar}%
}
\def\MyTOC@paragraph@star#1{%
  \MyTOC@orig@paragraph*{#1}%
}
\def\MyTOC@paragraph@nostar[#1]#2{%
  \MyTOC@orig@paragraph[{#1}]{#2}%
  % Si no numeras los \paragraph, cambia \theparagraph por texto plano sin \numberline
  \addcontentsline{stc}{paragraph}{\protect\numberline{\theparagraph}#2}%
}

% --- Impresor del índice propio (lee el .stc) ---
\newcommand{\MyTOC}{%
  \par\bigskip
  {\centering\bfseries\Large Índice de contenido\par}%
  \medskip
  \begingroup
    \parindent=0pt\parskip=2pt
    \@starttoc{stc}%
  \endgroup
}

\makeatother
% ===================== fin ToC propio =====================


%%%%%%%%%%%%%%


% --- Datos del tema ---
\title{Tema 3.A.42 -- Teorías de los ciclos económicos. \\
\large Ciclos nominales y reales. El ciclo financiero y sus interrelaciones con el ciclo real}
\author{Marcelo Ortega}
\date{Febrero 2026}
\newcommand{\TemaCode}{3A42}

\oldtitle{Tema 3.A.41. Teorías de los ciclos económicos: ciclos nominales y reales}
\changerationale{
    Los ciclos financieros han adquirido cada vez más relevancia, primero en el mundo académico, y luego en el de la política económica. En cualquier caso, se trataría de una extensión, no de una cuestión central en el tema y por tanto se recomienda su inclusión en el tema como tal (entre las numerosas críticas y ampliaciones al modelo de la NEK).
}

\begin{document}


\maketitle
\changebox

\newpage

% --- Página de resumen y modificaciones ---
\puntosclave{ %% Escribir puntos clave Apuntes CECO
    }
\vspace{1cm}

\modificaciones{
    \textcolor{red}{\textbf{OJO!}} Es necesario hacer especial énfasis en que desde una perspectiva histórica, la Teoría de Ciclos venía estrechamente ligada al desarrollo de modelos o teorías a través de las cuales poder \textbf{predecir} el ciclo. Es decir, en base a la evidencia empírica, poder establecer una serie de patrones estables de carácter longitudinal con el fin de proveer de información al policy-maker para anticiparse al origen del ciclo y la duración de este. 

    No obstante, a partir de los desarrollos de la RBC/TCNEK, el marco teórico se distancia de la idea de \textbf{anticipación}, \textbf{predictibilidad} y análisis longitudinal y se enfoca en buscar las \textbf{consecuencias} y \textbf{causas} del ciclo, de manera que la Teoría de Ciclos se centra en estudiar la \textbf{amplitud} e \textbf{implicaciones} del ciclo. Convergiendo, de esta forma, con el marco teórico de los modelos de OA-DA y el análisis de las perturbaciones y sus implicaciones de Política Económica.

    \textcolor{red}{\textbf{OJO!}} Algunos autores post-keynesianos mantienen como elemento central de su Teoría de Ciclos la \textbf{predictibilidad} y el análisis \textbf{longitudinal}. Por ejemplo, es característico el Modelo de Ciclos financieros de MINSKY. Podría ser interesante introducirlo por estas razones.
    

    \textcolor{red}{Recomendaciones Víctor}: Introducción —— 1.35

    Marco conceptual —— 5.40
    
    - Pondría grafico pizarra mientras explicas 
    - Algo más corto 
    
    Ciclos nominales —— 10
    
    - Curva Oferta Lucas, falta parte tendencial, poner una más simple. 
    - Dos problemas fundamentales con este modelo de ciclos? 
    
    RBC —— 13.32
    
    - Origen y transmisión de los ciclos en este caso? 
    - Poner la condición intertemporal 
    - Hay persistencia? Cómo se puede introducir? 
    - Qué es la elasticidad de sustitución intertemporal? 
    
    Ciclos NEK ——— 20.20 
    
    - Explicaria un poco más cómo has obtenido las funciones IS y NKCP
    - Qué es lo que explica que haya ciclos para la NEK? 
    - Algo más largo 
    
    Ciclos financieros —— 25.55 
    
    - Pondría gráfico 
    
    Hechos estilizados —— 28 
    
    - Los pondría en el primer epígráfe
}
    
    

\newpage
\MyTOC
\newpage

\mylistofequations
\clearpage

\MyListOfFigures
\clearpage


\pagenumbering{arabic}
\setcounter{page}{1}


\section{Introducción}\label{intro}


\subsection{Contextualización}\label{context}


\subsubsection{Relevancia}\label{importance}


\subsubsection{Historia}\label{history}



\subsubsection{Problemática}\label{problem} 




\subsection{Estructura de la exposición}\label{structure}






\clearpage
\section{Marco conceptual: ¿Qué es el Ciclo económico?}
\puntosclave{
    La actividad económica agregada ha estado desde siempre sujeta a fluctuaciones relevantes y las distintas escuelas de pensamiento económico han ido refinando progresivamente su análisis. Se puede distinguir dos grandes etapas: una previa a los modelos microfundamentados y la etapa más reciente basada en este tipo de modelos.
    
    En la etapa anterior a la microfundamentación del ciclo, se puede distinguir un primer período de “teoría de las crisis” o pre-ciclos (ausencia de una visión recurrente del ciclo), un segundo período de “edad de oro” de la teoría del ciclo (profusión de muy variadas teorías) que sienta las bases del análisis empírico del ciclo (Juglar, Mitchell…) y la aproximación estocástica al ciclo (Frisch y Slutsky), y un tercer periodo de desarrollo de teorías keynesianas del ciclo (Oxbridge), basado ya en un marco teórico macroeconómico propiamente dicho.

    La etapa de modelos de ciclo microfundamentados surge, como en el conjunto de la  macroeconomía, tras la crítica de Lucas y la revolución de las expectativas racionales en los años 70. Tras un primer intento con los ciclos nominales de Phelps-Lucas, se impondrán los modelos de ciclos reales (real business cycle o RBC), modelos de equilibrio general dinámico y estocástico (DSGE) en el marco (y con los supuestos) de la Nueva macroeconomía clásica. Pero su incapacidad para explicar satisfactoriamente algunas características de los ciclos llevó a que paulatinamente fueran sustituidos por los modelos neokeynesianos, que básicamente adaptan el modelo RBC a nuevos supuestos (competencia imperfecta, rigideces ...) y consiguen un mejor ajuste del modelo a la realidad observada.
}

\subsection{HEP: Análisis de los Ciclos Económicos} 
La existencia de crisis económicas caracterizadas por depresiones industriales ya se mencionaba en las obras de los mercantilistas (en particular, PETTY), los fisiócratas e incluso de SMITH. 

No obstante, es habitual identificar la crisis británica de 1825 como la primera crisis económica general de importancia que captó la atención de los economistas (a veces se usa esta fecha de 1825 como hito de inicio del capitalismo industrial moderno). Aunque el tema de las crisis se debatía desde hacía tiempo, el paso a la idea de que estas crisis son periódicas (es decir, algo recurrente o cíclico más allá de un incidente singular) llevó más tiempo en asentarse.

La evolución histórica del análisis de los ciclos puede dividirse, a grandes rasgos, en cuatro etapas:\footnote{Una versión extendida de estas etapas puede consultarse en los anexos.
}

\subsubsection{Teoría de Crisis: s.XVIII - 1860's}
La etapa de la Teoría de Crisis (mediados del s. XVIII - década de 1860), coincidiendo con el desarrollo de la escuela clásica, es una etapa pre-ciclos en la que se exploran las crisis o depresiones económicas como hechos aislados y sin un carácter recurrente (no hay ciclo strictu sensu). 

Estas teorías tratan en general de explicar las crisis económicas de manera aislada, como un suceso singular que sacude o perturba a las economías con carácter más o menos transitorio. No obstante, un minoritario grupo de autores consideraron la posibilidad de una crisis prolongada, como MALTHUS en la denominada controversia del \textit{general glut} (plétora generalizada) o MARX. 

En cualquier caso, su repetición en el tiempo (su carácter recurrente o cíclico) se veía como consecuencia de la repetición de los factores que las desencadenaban. 

No extraña esta falta de preocupación por el ciclo strictu sensu: el sistema capitalista moderno estaba en sus primeros años por lo que carecía de la historia suficiente como para detectar un patrón cíclico recurrente, a lo que se unía la falta tanto de una teoría económica suficientemente sofisticada que ayudase a comprender los fenómenos económicos (especialmente los macroeconómicos) como de un sistema de registro estadístico estable y fiable que permitiese conocer con rigor los hechos económicos objeto de estudio. 

\subsubsection{Teoría de los Ciclos de Negocios: JUGLAR (1862)}
La edad de oro de la teoría de los ciclos de negocios (década de 1860 – década de 1930), en la que se explora la existencia de ciclos y florecen muy variadas teorías de los ciclos, en parte debido a la amplia libertad que daba la ausencia de un marco o teoría macroeconómica propiamente dicha. Es gracias al estudio seminal de JUGLAR (1862) que los ciclos de negocios y su periodicidad quedaron establecidos como un tema de investigación. 

La teoría de los ciclos de negocios era considerada una cosa aparte, enraizadas más en la heterodoxia y los márgenes de la corriente neoclásica dominante, es decir, en las escuelas (hoy contempladas como) heterodoxas como la historicista o la institucionalista. Por tanto, durante gran parte de este período, la Teoría Económica y la Teoría de Ciclos de Negocios se movieron en planos separados y paralelos, sin apenas contacto y a veces incluso con hostilidad entre ellas.\footnote{%
    Cuando Adolph Lowe se preguntaba en un famoso artículo de 1926 sobre la cuestión \textit{How is business cycle theory even possible at all?}, la respuesta aún no estaba clara.
} 

En consecuencia, durante esta etapa se aprecia un enfoque dual:
\begin{la}
    \item \textbf{Enfoque descriptivo}. Por una parte, un enfoque meramente descriptivo o empírico del ciclo económico (W. C. MITCHELL); y, 
    \item \textbf{Enfoque normativo}. Por otro, un conjunto de teorías que resultaban insatisfactorias por incompletas: (i) bien porque vinculan el ciclo económico a otro tipo de ciclos (climáticos, de crédito, expectativas, etc.) que quedaban inexplicados; o, (ii) bien porque lo atribuían a fenómenos puramente estocásticos (FRISCH-SLUTSKY).\footnote{%
        Muchas de estas explicaciones eran conjeturas ad hoc construidas alrededor de hechos empíricos más que derivadas de fundamentales subyacentes. Se encuentran así teorías basadas en el clima (JEVONS, ÅKERMAN), teorías de la sobreinversión (JUGLAR, SPIETHOFF, ROBERTSON), teorías psicológicas (AFTALION, CLARK, PIGOU), teorías monetarias (HAWTREY, HAYEK), teorías del subconsumo (HOBSON, FOSTER y CATCHINGS) y teorías basadas en shocks aleatorios (FRISCH, SLUTSKY).
    }
\end{la}

Curiosamente, conviene destacar que la teoría del ciclo actual es fuertemente tributaria de ambos enfoques: de los esfuerzos empíricos MITCHELL y el NBER, de la escuela institucionalista norteamericana, por determinar las características y regularidades del ciclo, y del enfoque teórico FRISCH-SLUTSKY, de atribuir al ciclo un origen exógeno y aleatorio. Fue únicamente en la década de los años 1920's del pasado siglo cuando los economistas empezaron a tratar de encontrar la manera de conectar los ciclos de negocios con la teoría económica propiamente dicha. Aunque estas eran aún tentativas de conexión.

\subsubsection{Etapa keynesiana: Teoría Macroeconómica}
La etapa keynesiana (década de 1930 – década de 1970), en que la Teoría del Ciclo posee por primera vez un marco analítico general y estable: la teoría macroeconómica. 

El programa de investigación de los seguidores de Keynes (llamado Oxbridge por incorporar a autores tanto de Oxford como de Cambridge) consistió así en adaptar la teoría de aquel para conseguir explicar el ciclo económico. Se pueden dividir los diferentes esfuerzos teóricos de la escuela keynesiana en dos:
\begin{la}
    \item Las teorías del \textbf{multiplicador-acelerador} (SAMUELSON, HICKS, PASINETTI); y,
    \item Las teorías de \textbf{ciclos endógenos} (KALECKI, KALDOR, GOODWIN).
\end{la} 

Un común denominador de estas teorías es, paradójicamente, su limitada inspiración keynesiana. Esto queda patente en dos aspectos fundamentales: (i) frente a la visión de KEYNES del ciclo causado por la dinámica de las expectativas empresariales o animal spirits, que afectaría a la eficiencia marginal del capital y de ahí al multiplicador de la renta, estas teorías son totalmente deterministas; y, (ii) comparten un deliberado desinterés por factores monetarios o financieros como causas del crecimiento o los ciclos. Todos ellos contemplaron al mercado como creador del crecimiento y de los ciclos y desarrollaron modelos que trataban, paradójicamente para ser keynesianos, de generar estos fenómenos dinámicos mayoritariamente en sistemas reales, alejados de consideraciones monetarias o nominales.

Las recesiones eran periodos con una baja utilización de recursos que podrían ser estimulados por medio de políticas dirigidas a expandir la demanda agregada. Otro de los aspectos conceptuales de las teorías RBC es la importancia de los shocks tecnológicos (calibrados para ajustar el residuo de Solow o productividad total de los factores o PTF) como fuente de las fluctuaciones económicas, y del papel limitado de los factores monetarios para explicar las fluctuaciones económicas.

\subsection{Objeto de estudio: Implicaciones de Política Económica}
Es gracias a este desarrollo histórico, a través del cual la Teoría de Ciclos ha ido encontrando un lugar propio dentro de la gran cantidad de ramificaciones que caracteriza a la Cicencia Económica, por lo que podemos definir, con cierto consenso, algunos de los elementos/características esenciales de la Teoría de Ciclos:\footnote{%
    Como se ha dicho, en los años 30, los econometristas A. F. BURNS y W. C. MITCHELL empezaron a documentar la existencia de un conjunto de regularidades empíricas del ciclo de negocios. Esta investigación culminó en su tratado de 1946, \textit{Measuring Business Cycles}, en la que aparece una definición del ciclo de negocios que ha devenido clásica:
\begin{cita}{A. F. BURNS y W. C. MITCHELL}{Measuring Business Cycles}{1946}
    Los ciclos económicos son fluctuaciones en la actividad agregada de las economías de mercado: un ciclo consiste en expansiones que ocurren al mismo tiempo en muchas actividades económicas, seguido de recesiones o contracciones igualmente generalizadas, y reactivaciones que conducen a la fase de expansión del siguiente ciclo; esta secuencia de cambios es recurrente, pero no periódica; en cuanto a su duración, los ciclos económicos varían desde más de un año hasta diez o doce años; no son divisibles en ciclos más cortos de duración exacta.
\end{cita}

\textbf{NOTA.-} Se sigue el análisis de SØRENSEN y WHITTA-JACOBSEN (2010). Se aconseja consultar esta fuente para profundizar: lo aquí presentado es una visión superficial.
} 
\begin{lnum}
    \item \textbf{Perturbaciones sobre la economía en su conjunto}. Afectan a la actividad agregada (no solo PIB, sino movimiento simultáneo de gran número de variables: empleo, tipos de interés, precios, etc.); 
    \item \textbf{Economía de mercado}. Tiene lugar en economías de mercado (países que organizan sus economías a través de “empresas de negocios”, pero no en las antiguas economías socialistas); 
    \item \textbf{Patrón: Pico-Valle}. Expansiones seguidas de contracciones: se sigue un patrón pico – valle;
    \item \textbf{Recurrentes pero no periódicos}. Es decir, carecen de patrón regular sencillo, pues ni la ocurrencia ni la amplitud son perfectamente regulares y, por tanto, tampoco son fácilmente predecibles. Se ha desistido de explicar los ciclos como combinación de ciclos deterministas de diferente periodicidad (KITCHIN de 30-40 meses, JUGLAR de 3-5 años, KUZNETS de 16-22 años, KONDRATIEFF de 50 años); y,
    \item \textbf{Persistencia}. Duración superior al año, lo que implica que: (i) las fluctuaciones inferiores al año, de carácter estacional, no tienen la consideración de ciclos; (ii) existe persistencia de unos supuestos shocks puntuales iniciadores del ciclo (elevada autocorrelación serial).
\end{lnum}

Por tanto, para conocer las características esenciales de las fluctuaciones cíclicas es esencial: (i) poder delimitar temporalmente el ciclo; y, (ii) obtener las propias series temporales de las variables de estudio.

\subsubsection{Delimitación temporal: Patrón} 
BURNS y MITCHELL ofrecieron por primera vez una datación de los ciclos estadounidenses, mediante la identificación de los puntos de cambio de tendencia (turning points) de la actividad económica, de modo que el ciclo se puede medir de valle a valle. 

Desde entonces, el NBER estadounidense mediante su comité de expertos ha realizado su seguimiento para la economía estadounidense y define una recesión como un periodo de declive significativo de la producción, la renta, el empleo y el comercio, y marcada por una contracción generalizada en muchos sectores de la economía y que dura desde 6 meses a 1 año. Esta longitud de valle a valle ha variado mucho en los distintos períodos históricos (desde 4 trimestres en 1980-81 a casi 10 años en 1960-70 y 1991-2000), siendo las fases expansivas más largas que las contractivas. La principal conclusión que se extrae es que las recesiones y expansiones difieren en su tamaño, duración y frecuencia de ocurrencia, por lo que se ha huido de interpretaciones basadas en combinaciones de ciclos determinísticos de distinta longitud (Kitchin, Juglar, Kuznets, Kondratieff), y la visión que prevalece es que la economía está expuesta a perturbaciones de distintos orígenes y tamaños, que se producen a intervalos aleatorios, y que se propagan por toda la economía. Las discrepancias entre las principales escuelas de pensamiento residen en las hipótesis sobre cuáles son los orígenes de las perturbaciones y los mecanismos de propagación (Romer, 2018).

La identificación de valles y picos se basa en unas sencillas reglas:
\begin{lnum}
    \item Un valle debe ser seguido por un pico, y un pico debe ser seguido por un valle; 
    \item La fase de expansión, así como la de contracción debe durar como mínimo 2 trimestres: este criterio implica que requerimos cierto grado de persistencia en el movimiento de la actividad económica; y,
    \item Un ciclo completo de negocios debe contener al menos 5 trimestres: este criterio implica que las fluctuaciones que duran menos de un año no se consideran ciclos de negocios.
\end{lnum}

\subsubsection{Medición de los ciclos: ¿errores de medida?}
Los valores observados de las variables macroeconómicas de estudio son la combinación de tres componentes: tendencial, cíclico y estacional. En la realidad sólo observamos su suma, por lo que se requieren técnicas para separarlas.

\paragraph{Componente estacional: TRAMO}
El componente estacional se suele eliminar en primer lugar de modo que las series queden limpias de las fluctuaciones periódicas que se producen dentro del año. Este paso suele realizarse mediante técnicas estadísticas avanzadas.\footnote{%
    Una de las herramientas más utilizadas internacionalmente son los programas TRAMO (Time series Regression with ARIMA noise, Missing values and Outliers) y SEATS (Signal Extraction in ARIMA Time Series), desarrollados por Gómez y Maravall, del Banco de España.
}

\paragraph{Componente tendencial y cíclico: ¿?}
De la serie desestacionalizada se separaría el componente cíclico ($y_t^c$) del tendencial ($y_t^g$), en logaritmos:

\eqblock{
\begin{aligned}
    y_t^{\text{\tiny observado}}\quad \overset{\text{\scriptsize $y_t^e\not\in y_t$}}{\underset{\text{\scriptsize TRAMO}}{\Longrightarrow}}\quad \boxed{y_t=y_t^g+\underset{\text{\scriptsize objeto}}{y_t^c}}
\end{aligned}
}{Variable observada: Descomposición. Marco conceptual: objeto de estudio}

Como no se conoce con certeza el funcionamiento de las economías (de ahí el uso de diferentes modelos para captar tal funcionamiento y sobre los que discutiremos en las secciones siguientes), no existe un único método correcto para separar la tendencia del ciclo. 

Una tendencia lineal puede descartarse, pues implica aceptar que las economías están siempre en un hipotético estado estacionario (creciendo a una tasa constante). En efecto, la teoría moderna del crecimiento endógeno sugiere que la tasa de progreso técnico puede variar con la actividad de innovación endógena de las empresas (economía de las ideas) o con la acumulación de capital humano, por lo que esa tendencia no será lineal, sino que presentará sus propias fluctuaciones.\footnote{%
    Para ver los diferentes métodos disponibles para la descomposición de la variable observada: [\authorfont{Ver Anexo \ref{anx:filtros_descomposicion}}]
}

\clearpage
\section{Teoría del Ciclo Real: Equilibrio General Competitivo}
\puntosclave{
    Los modelos de ciclo real, así como los posteriores neokeynesianos, son modelos de crecimiento microfundamentados (Ramsey-Cass-Koopmans), de tipo walrasiano (sin fallos de mercado ni fricciones) y ampliados con: a) perturbaciones estocásticas, b) oferta de trabajo endógena, y c) mayor detalle en las formas funcionales (utilidad, producción…). 
    
    Se trata de modelos de equilibrio general (con consumidores y empresas optimizadoras y habitualmente sin sector público), dinámico (horizonte infinito) y estocástico (perturbaciones aleatorias).
    
    Como resultados principales, los modelos de ciclo real consideran que el origen del ciclo es exógeno, real (habitualmente en la productividad total de los factores (PTF)) y aleatorio. La propagación se garantiza con su naturaleza de equilibrio general. La persistencia se buscó inicialmente a través de la dinámica del capital o de la elevada elasticidad de sustitución intertemporal del trabajo, pero solo se consigue incorporando como supuesto perturbaciones autorregresivas. Los modelos de ciclo real presentan una elevada complejidad técnica que dificulta su resolución y, más importante, su comprensión.
    
    Como valoración, la evidencia abrumadora es que los modelos de ciclo real no consiguieron capturar las características clave de las fluctuaciones observadas. En efecto, implican neutralidad monetaria, requieren perturbaciones tecnológicas que a duras penas se observan en la realidad, se basan en fundamentos microeconómicos sin apenas respaldo empírico y arrojan dinámicas de las variables que difieren significativamente de las observadas en la realidad. Estas deficiencias dieron lugar a una abundante literatura de extensiones “walrasianas” de los modelos iniciales hasta finales de los años 90.
}

Como hemos visto, el increíle avance en términos estadísticos experimentado a lo largo de la segunda mitad del siglo XX hizo que la Teoría Económica se viese suparada a la hora de proporcionar los fundamentos económicos (teóricos) que diesen sentido a las obervaciones. 

Más aún, la incapacidad de dar recomendaciones de Política Económica acordes con la situación que se venía experimentado desde finales de los años 60's y que se vió especialmente agravada con el inicio de los años 70's y las fuertes perturbaciones macroeconómicas derivadas de las sucesivas crisis del petróleo; sumió al mundo académico en una crisis que derivó en la revolución metodológica de la NMC: la crítica de LUCAS. Esta crítica al uso de modelos macroeconométricos realizados hasta ese momento así como la revolución de la HER hizo que volvieran a utilizarse modelos de Equilibrio General micro-fundamentados para estudiar los ciclos económicos. 

Dió así inició lo que se conoce como la Teoría del Ciclo Real o RBC, que expondremos a continuación.

\textbf{¿Qué queremos decir?} - Estos autores defienden: (i) la coherencia entre las Teorías del Ciclo y del crecimiento económico (por lo que los modelos de ciclo son en esencia modelos de crecimiento con ciertas adaptaciones); y, (ii) la naturaleza eficiente de éstos. Es decir, consideran que se trata de respuestas de ajuste de la economía ante shocks reales (sobre todo shocks tecnológicos, pero también a shocks en preferencias identificados como dedemanda), lo que contradecía la visión keynesiana. No obstante, también veremos que a pesar de su elegantísimo marco conceptual: (i) no son capaces de resistir la contrastación empírica; y, (ii) consideran que los shocks tienen, principalmente, naturaleza tecnológica (lo cual no es consistente con las observaciones).

\subsection{Ciclos nominales: Curva de Oferta de LUCAS}
La primera aportación de esta linea de investigación es la conocida como Curva de Oferta de LUCAS. Se trata de un modelo simple de ciclos nominales de información imperfecta, que se utilizó para demostrar la ineficiencia derivada de los ajustes de Política Monetaria que se llevasen a cabo para contrarrestar el ciclo. 

LUCAS considera un contexto de equilibrio general dinámica, vaciamiento de mercados continuo y agentes con conducta racional (maximizadores de su función objetivo). De este forma, podemos representar la Curva de Oferta de LUCAS (sin entrar en su derivación analítica) como:

\eqblock{
\begin{aligned}
    &\text{Curva de Oferta de LUCAS}:&&\boxed{y_t^s=\beta(1-\theta)(p_t-\mathbb E_t[p_t])} \quad \beta\equiv\varepsilon_{S}\\[2em]
\end{aligned}
}{Curva de Oferta de LUCAS: Perturbaciones reales y nominales. Teoría de los Ciclos Reales: Ciclos nominales}

¿Qué quiero expresar LUCAS con esta especificación? Como vemos, la variación dinámica de la oferta agregada depende de dos factores cruciales: (1) el parámetro $\theta$; y, (ii) la brecha entre el nivel de precios agregado observado y las expectativas de los agentes respecto a al nivel de precios. En consecuencia, 

\eqblock{
\begin{aligned}
    &\underset{\text{\scriptsize (se muestran los casos polares)}}{\text{Si observa $p_t-\mathbb E_t[p_t]>0$}}:
    \left\{\begin{aligned}
        &(1-\theta)=1\quad\Longrightarrow\quad\underset{\text{\scriptsize schock real}}{y_t^s>0}\\[1em]
        &(1-\theta)=0\quad\Longrightarrow\quad \underset{\text{\scriptsize schock nominal}}{y_t^s=0}
    \end{aligned}\right.\\[1em]
    &\text{donde,}
    \underset{\substack{\\[1em]\\\text{\scriptsize $\Longrightarrow$ Límites de la Política Monetaria estabilizadora: Confianza de los agentes}}}{\left\{\begin{aligned}
        &(1-\theta)\to1\quad\iff\quad p_t -\mathbb E_t[p_t\;|\;p(z)_t]>0\quad:\text{No anticipado}\\[0.5em]
        &(1-\theta)\to0\quad\iff\quad p_t -\mathbb E_t[p_t\;|\;p(z)_t]=0\quad:\text{Anticipado}\\[0.5em]
    \end{aligned}\right.}
\end{aligned}
}{Curva de Oferta de LUCAS: Impicaciones de Política Monetaria. Teoría de los Ciclos Reales: Ciclos Nominales}

Por tanto, LUCAS demuestra (analíticamente) que la brecha entre el nivel de precios observado y las expectativas formadas por los agentes sobre éste, tendrá consecuencias reales si y solo si los agentes otorgan cierta credibilidad a la Autoridad Monetaria. 

\subsubsection{Implicaciones de Política Económica}
De esta forma, las implicaciones de Política Económica son profundas y variadas, de entre las cuales cabe destacar dos:
\begin{lnum}
    \item \textbf{Ineficacia de la Política Monetaria}. La explotación recurrente del trade off entre inflación y desempleo conlleva que la Política Monetaria devenga absolutamente ineficaz. Esto es así porque: (i) desde una perspectiva analítica, los agentes reformularán los pesos estadisticos que otorgan a la varianza en el precio observado local y el nivel agregado de precios; y, (ii) intuitivamente, la pérdida de credibilida de la Autoridad Monetaria llevará a los agentes a una situación límite de superneutralidad del dinero (absorber todo estímulo monetario a través del incremento de los precios);
    \item \textbf{Shocks de naturaleza real}. En un contexto de plena credibilidad de la Autoridad Monetaria, esta será capaz de influir en la economía y generar shocks nominales debido a la inflación no anticipada. 
\end{lnum}

Pero más interesenta aún, y lo que resulta ser el elemento esencial de la Teoría de Ciclos Reales: las únicas perturbaciones económicas (no oportunistas) que pueden originar los ciclos son de naturaleza real.

\subsection{Ciclos Reales: Modelo canónico}
Lo que nos conduce al verdadero modelo canónico de la Teoría de Ciclos Reales: el Real Bussiness Cycle propuesto por KYDLAND y PRESCOTT.

Estos autores sobre la base de que el origen último del ciclo es de carácter real, homogeneizan estas perturbaciones bajo el paraguas de \textbf{perturbación tecnológicas} (en ausencia de intervenciones no anticipadas del policy maker que generen fluctuaciones nominales $\to$ reales).

Por tanto, empezamos definiendo la dinámica de la tecnología capaz de captar esto:

\eqblock{
\begin{aligned}
    \boxed{a_t=a_t^g+a_t^c}\qquad \text{donde}\quad \boxed{a_t^c=\rho_aa_{t-1}^c+u_t^a}\qquad 
    \left\{\begin{aligned}
        &u_t^a\sim N(0,\sigma_a)&&\equiv\text{Perturbación estocástica}\\[0.5em]
        &\rho_a&&\equiv\underset{\text{\scriptsize resultados muy sensibles a este parámetro}}{\text{Parámetro de persistencia}}
    \end{aligned}\right.
\end{aligned}
}{Tecnología: Componente lineal y estocástico. Teoría de Ciclos Reales: KYDLAND y PRESCOTT}

Por tanto, este tipo de shock o perturbación se introduce como componente aleatorio (esto justifica el calificativo de estocástico del modelo), formado a su vez por un parámetro de persistencia y el término de error que genera la perturbación estocástica.\footnote{%
    Aunque se pueden encontrar otras fuentes de perturbación real, como las habidas en el gasto público o las preferencias, aquí nos centraremos en el caso más puro o fiel al origen de estos modelos. Claro está, con posterioridad dentro de esta corriente se añadirían “complicaciones” o refinamientos que aumentarían, aunque de manera limitada, la capacidad explicativa del modelo.
} No obstante, cabe destacar que la persistencia no es un resultado del modelo, sino que se impone ex ante como supuesto, haciendo en ocasiones que $\rho_a$ sea el parámetro al que los resultados del modelo son más sensibles. 

Ahora que hemos definido la fuente de perturbación, veamos el resto de componentes del modelo.

\subsubsection{Optimización de los hogares}
Los consumidores u hogares maximizan su utilidad, dependiente del consumo y del trabajo. Asumimos que la FUI es separable y aditiva y se plantea en un horizonte de dos periodos. De esta forma, el probla de maximización del consumidor puede representarse como:

\eqblock{
\begin{aligned}
   &\boxed{\begin{aligned}
        &\max_{\{c_1, c_2, n_1, n_2\}}{U(0)}=u(c_1, n_1)+\beta(c_2, n_2)\\[1em]
        &\text{s.a.}\quad c_1+\frac{1}{1+r}c_2=w_1n_1+\frac{1}{1+r}w_2n_2
    \end{aligned}}\to\text{Resolvemos por}\; \mathcal{L}\;\text{y obtenemos CPO's}\\[1em]
    &\boxed{\frac{u'(n_1)}{u'(n_2)}=\frac{w_1}{w_2}(1+r)\beta}\to \text{Condición intertemporal de Ocio}
\end{aligned}
}{Problema de optimización del consumo: Decisión intertemporal consumo-trabajo. Teoría de Ciclos Reales}

La restricción de recursosse obtiene de: (i) trabajo a cambio de un salario real (riqueza humana); y, (ii) riqueza no humana (capital/ahorro) que genera rentabilidad. 

\paragraph{Perturbación tecnológica: Efectos sobre la condición de EULER}
En general se pueden contemplar los efectos del shock operando a través de efectos riqueza y de sustitución intertemporal:

\eqblock{
\begin{aligned}
    a_t^c>0\quad\Longrightarrow\quad
    \left\{\begin{aligned}
        &(1)\text{Efecto sustitución}:&&
        \left\{\begin{aligned}
            &\underset{\text{\scriptsize aumenta la productividad}}{F'(n_t)>F'(n_{t-1})}\\
            &\underset{\text{\scriptsize escasez relativa de capital}}{\frac{F'(n_t)}{F'(k_t)}>1}
        \end{aligned}\right\}\Longrightarrow
        \left\{\begin{aligned}
            &{w_t>0}\Rightarrow n_t>0\\
            &{r_t>r_{t-1}}\Rightarrow s_t>0
        \end{aligned}\right\}\\[1em]
        &(2)\text{Efecto riqueza}:&&\sum_{t=0}^T\beta^ts_t^t>\sum_{t=0}^T\beta^ts_{t}^{t-1}\Longrightarrow c_t>0
    \end{aligned}\right.
\end{aligned}
}{Efectos de un shock tecnológico positivo: Efecto riqueza y efecto sustitución. Teoría del Ciclo Real: KYDLAND y PRESCOTT}

Un shock positivo transitorio de productividad implica que la economía será más productiva por un tiempo. Por tanto: (i) por un lado, la riqueza intertemporal de los agentes será mayor, lo que les incitaría a aumentar su consumo presente (efecto riqueza); y, (ii) por otro lado, aumentarán la oferta de trabajo y ahorrarán más (efecto sustitución intertemporal) porque la mayor productividad es transitoria, luego es sensato aprovecharla mientras dure y trabajar más, y el stock de capital es bajo con relación a la tecnología (dado el aumento de PTF), así que la productividad marginal del capital (y con ello la remuneración del ahorro) es especialmente elevada, lo que lleva a ahorras más. 

Por lo tanto, el efecto sustitución intertemporal debe dominar al efecto riqueza para que la oferta de trabajo y el ahorro crezcan en el corto plazo. 

\subsubsection{Optimización de las empresas}
Por su parte, las empresas maximizan su beneficio operando en un entorno de competencia perfecta y cuya tecnología se resume en una función clave, la función de producción agregada (FPA), que expresa la cantidad total de producto (homogéneo, para toda la economía) en función de las cantidades de los factores primarios, trabajo y capital, que demanda a los hogares. Suponiendo una FPA de tecnología Cobb-Douglas:

\eqblock{
\begin{aligned}
    &\max_{K_t,L_t} \Pi(0)=\int_0^\infty e^{-rt}[P_xF(K_t,L_t)-r_tK_t-w_tL_t]\mathrm dt\\[0.5em]
    &\text{s.a.}\quad \dot K_t=K_t+(1-\delta)K_{t-1}
\end{aligned}
}{Mercados perfectos: No hay conexión intertemporal. Teoría de Ciclos Reales}

Para la empresa representativa, al igual que en muchos modelos de crecimiento económico, típicamente su problema de optimización no suele ser genuinamente intertemporal, pues ninguna de sus decisiones presentes sobre los factores limitará dichas decisiones en el futo debido a que el fondo de capital puede ajustarse perfectamente encada período y carecen de una verdadera restricción de recursos al tener acceso a una financiación sin restricciones (mercados de capitales perfectos). Un problema prototípico (caso determinístico) de una empresa representativa que maximiza la senda de beneficios futuros (ingresos por ventas de la producción menos costes de inversión y de trabajo en cada período), con un único bien, tanto de inversión como de consumo, cuyo precio se normaliza a 1, sería:


La PTF contine una tendencia lineal21, ����, y un componente aleatorio, �� ̃ ��. Esto implica que todas las variables endógenas futuras, tanto precios como cantidades, sean estocásticas, lo que obliga a aplicar la teoría de la utilidad esperada. De este modo, los consumidores maximizan la esperanza de sus utilidades futuras, formando dichas expectativas racionalmente (hipótesis de expectativas racionales). Sin embargo, las empresas se ven exentas de esta complicación, pues como se ha indicado su problema no es genuinamente intertemporal.

Dado que una característica que se exige a la teoría es la persistencia, en un primer momento el instrumento para conseguirla fue el capital, de modo que el shock inicial provoque una “sobrerreacción” en forma de inversión/acumulación de capital que perdure múltiples períodos a través de la ecuación de movimiento del fondo de capital (ya incorporada en la función objetivo anterior): ����+1 = ���� + ���� − ������.

Respecto a la dinámica del capital, se llega a demostrar que no aporta apenas capacidad explicativa al modelo,\footnote{%
    Como demostraron célebremente Cogley y Nason (1995), la respuesta de la inversión y del stock de capital a los shocks de productividad contribuye en realidad muy poco a la dinámica de estos modelos. Hasta el punto de que la siguiente familia de modelos, los neokeynesianos, pueden prescindir totalmente de la variable capital.
} lo que lleva a introducir shocks que siguen procesos autorregresivos.\footnote{%
    Lo que significa que el valor del shock en un período dependa de los valores que tomó en períodos anteriores, luego se garantiza/impone que, si se produce un shock en un período determinado, en los siguientes persistirá.
}







Un aspecto técnico, pero de especial relevancia para valorar el modelo RBC (y cualquier modelo EGDE), es la resolución de este dada su elevada complejidad. Aunque existen diferentes opciones para solucionar modelos DSGE, el común denominador es una creciente falta de transparencia: se puede llegar a conocer, mediante funciones de impulso-respuesta, la senda exacta que seguirán las variables endógenas ante determinados shocks o políticas, pero no se podrán conocer los mecanismos lógicos que expliquen racionalmente por qué esas variables relevantes responden de una determinada manera a dichas perturbaciones.


Y respecto a la elasticidad de sustitución intertemporal (ESI), siguiendo a Romer (2018), 

Suponiendo una función de producción agregada (FPA) de tecnología Cobb-Douglas , la tecnología contiene un componente estocástico. 

Se requiere que se produzcan variaciones en el empleo. Frente al modelo base de crecimiento Ramsey-Cass-Koopmans, en que todo el trabajo se oferta inelásticamente (pues no afecta a la utilidad), se necesita incorporar al consumo (����) el ocio o el trabajo (����) en la función de utilidad, ��(����, ���� ). Esto permitirá obtener una oferta de trabajoendógena que responda al salario real.




\subsection{Implicaciones de Política Económica: Eficiencia del Ciclo}
En definitiva, la Teoría de Ciclos Reales es el primer modelo que permite explicar el ciclo de manera microfundamentada. Esto es, explicando el comportamiento agregado de la economía como consecuencia de los comportamientos microeconómicos de los agentes y, en consecuencia, la consideración de fundamentos económicos sólidos (o, como mínimo, de consenso): el origen del ciclo es exógeno, real y aleatorio, inicialmente tecnológico vía un shock transitorio de la PTF y la dinámica de ajuste observada responde al comportamiento de los agntes antes este shock.\footnote{%
    Posteriormente se amplió el abanico de perturbaciones reales a: (i) las variaciones en las compras públicas (Aiyagari, Christiano y Eichenbaum,1992; Baxter y King, 1993); (ii) noticias acerca de cambios futuros en la tecnología de la economía (Beaudry y Portier, 2004, 2006; Jaimovich y Rebelo, 2009; Alexopoulos, 2011); o, (iii) perturbaciones en la tecnología para producir bienes de inversión (Greenwood, Hercowitz y Huffman, 1988; Hornstein y Krusell, 1996).
} Además, el hecho de que este shock puntual se propague al conjunto de variables de la economía segarantiza con la naturaleza de Equilibrio General del modelo, que no es recursivo (no puede dividirse en bloques separados que se resuelven de manera independiente).\footnote{%
    En concreto, para el caso analizado:

El shock de PTF impacta directamente en las productividades marginales de trabajo y capital y, por tanto, altera las remuneraciones de ambos factores, afectando tanto a las decisiones básicas de empresas como de consumidores, que modificarán en consecuencia sus sendas de oferta y demanda de factores (trabajo e inversión) y de producción y consumo. Y la persistencia se obtiene o por la vía de la dinámica de la acumulación de capital o por la elasticidad de sustitución intertemporal (ESI) de los consumidores.
}

No obstante, por construcción, sus consecuencia/implicaciones en términos de Política Económico son desoladoras. En efecto, al utilizar un modelo walrasiano (sin ningún tipo de imperfección ni fricciones) y sin heterogeneidad de agentes (es de agente representativo): el ciclo económico es \textbf{eficiente} en sí mismo pues es el comportamiento de la economía se ajusta correctamente a todos los parámetros fundamentales que alteran el comportamiento dinámico de los agentes frente a una perturbación. En consecuencia, la asignación PARETO-Óptima alcanzada impide que la intervención del policy-maker encuentra ninguna justificación económica.

Por eso precisamente no incorpora Sector Público: dado que no hay fallos de mercado de ningún tipo, los ciclos son respuestas óptimas de los agentes a los shocks, por lo que intervenir para combatirlos sería ineficiente (el equilibrio competitivo alcanzado es en todo momento óptimo de PARETO). La recomendación de política económica es desaconsejar las políticas macroeconómicas estabilizadoras. 

\subsubsection{Limitaciones: Fundamentos económicos}
Dicho esto, estas consecuencias deben tomarse con muchísima cautela pues presenta grandes problemas que limitan su veracidad y uso. El modelo de Ciclo Real es la extensión más simple y natural del modelo de RAMSEY para incluir fluctuaciones. Y, aunque esta simpleza permite descubrir unos mecanismos esenciales de propagación y trasmisión, originando un modelo base o de referencia, esto tiene un precio: la evidencia abrumadora es que no consiguen capturar las características clave de las fluctuaciones observadas. Podemos distinguir tres limitaciones principales.

\paragraph{Shocks tecnológicos: ¿cuántos hay?}
El modelo RBC se basa crucialmente en grandes shocks tecnológicos agregados, de los que hay escasa evidencia. En efecto, es habitualmente difícil identificar innovaciones específicas asociadas a variaciones intertrimestrales de la PTF. 

Aún más importante, existe significativa evidencia de que las variaciones a corto plazo del residuo de SOLOW incluyen algo más que innovación tecnológica (Bernanke y Parkinson, 1991; Mankiw, 1989; Hall, 1988), por lo que el shock verdadero sería incluso inferior, lo que reduciría el poder explicativo del modelo. Y cuando se han identificado verdaderos shocks tecnológicos y sus efectos, se deduce que algunas variables (empleo) se comportan de forma opuesta a la predicha por el modelo RBC (Shea, 1998; Galí y Rabanal, 2004; Francis y Ramey, 2005; Fernald, 2007).

Otro de los aspectos conceptuales de las teorías RBC es la importancia de los shocks tecnológicos (calibrados para ajustar el residuo de Solow o productividad total de los factores o PTF) como fuente de las fluctuaciones económicas, y del papel limitado de los factores monetarios para explicar las fluctuaciones económicas.\footnote{%
    Es más, la mayoría de los modelos RBC se abstrae de modelizar un sector monetario o una Autoridad Monetaria. Incluso la introducción de una Autoridad Monetaria en modelos con mercados sin fricciones y actuando en competencia perfecta, como el de Cooley y Hansen (T-F. Cooley y G.P. Hansen, 1989, lnflation Tax in a Real Business Cycle Model, American Economic Review, 79(4), págs. 733-7481) (y otros) no fueron percibidos como que proporcionaran una estructura atractiva para el análisis de política monetaria, ya que estos modelos predecían neutralidad monetaria de la política monetaria sobre la actividad real.
} 

\paragraph{Calibración: Parámetros excesivos}
Dado que el objetivo es explicar con cierto detalle un gran número de macromagnitudesy sus características, esto obliga a especificar con más detalle las funciones del modelo (las funciones de utilidad y producción ya no pueden ser genéricas) para poder alcanzar una solución, lo que se traduce en una mayor cantidad de parámetros o variables exógenas del modelo. Y dado que se necesita replicar o simular el comportamiento observado de la economía para validar el modelo, la determinación de los valores correctos de esos parámetros se alcanza mediante ejercicios de calibración.\footnote{%
    La calibración de un modelo consiste en elegir el vector con los valores exactos de los parámetros exógenos del modelo. Para ello se recurre a evidencia empírica sobre esos parámetros, que se ajustará haciendo simulaciones retrospectivas con el modelo hasta que ofrezca resultados cercanos a los datos reales (en términos de varianzas y covarianzas de diferentes series). Véase apartado 5.8 de Romer (2018).
}

Los fundamentos microeconómicos del modelo parecen demasiado alejados de lo que se observa en la realidad (características no walrasianas de las economías actuales, así como una baja ESI, clave para que el modelo funcione). 


\clearpage
\section{Teoría Neokeynesiana del Ciclo}
\puntosclave{
    Los modelos neokeynesianos de ciclo son a su vez una extensión de los modelos RBC, fundamentalmente son modelo EGDE que incorporan competencia imperfecta, rigideces nominales y múltiples fuentes de perturbaciones (reales y nominales).

    Se trata igualmente de modelos de equilibrio general (con consumidores y empresas optimizadoras y ahora también con sector público en su rol de autoridad monetaria y a veces fiscal), dinámico (horizonte infinito) y estocástico (perturbaciones aleatorias de  diversa naturaleza). 

    Como resultados principales, los modelos neokeynesianos de ciclo también consideran que el origen del ciclo es exógeno y aleatorio, pero puede ser tanto real como nominal (monetario). La propagación se garantiza igualmente con su naturaleza de equilibrio general. Pese a las rigideces nominales, la persistencia también se debe imponer con perturbaciones autorregresivas. Los modelos de ciclo neokeynesianos presentan una mayor complejidad técnica que los modelos RBC, agravando los problemas de resolución y comprensión.
    
    Como valoración, los modelos neokeynesianos de ciclo suponen una mejora respecto a los modelos RBC, pues exhiben no neutralidad de la política monetaria, un mejor apoyo empírico microeconómicos en sus supuestos y un mejor ajuste a los datos observados. Sin embargo, también presenta debilidades, como la incapacidad para exhibir per se la suficiente persistencia de las perturbaciones o la conocida como forward guidance puzzle. Esto ha dado lugar a numerosas extensiones y refinamientos de los modelos iniciales.
    
    Dado que la teoría neokeynesiana abre la puerta a la intervención pública estabilizadora, es imprescindible preguntarse por su justificación, pues no es obvio cómo la inflación y las fluctuaciones en la producción afectan al bienestar. A fecha de hoy, nuestro conocimiento sobre los costes de la inflación y de las variaciones de la producción que deberían guiar esta política estabilizadora sigue siendo abrumadoramente limitado.
}

Desde un punto de vista metodológico, la teoría RBC estableció el uso de Modelos de Equilibrio General Dinámicos y Estocásticos como la herramienta central, no sólo en el análisis de ciclos, sino para todo el análisis macroeconómico.\footnote{%
    Así, las ecuaciones de comportamiento que describían las relaciones entre variables macroeconómicas utilizadas hasta entonces por la macroeconomía keynesiana, fueron sustituidas por las condiciones de primer orden de los problemas intertemporales a los que se enfrentaban los agentes económicos modelizados (hogares y empresas). Los supuestos ad hoc acerca de la formación de expectativas que se habían utilizado hasta ese momento (expectativas estáticas o adaptativas) dieron paso al supuesto de que los agentes formaban expectativas racionalmente, es decir, no cometían errores sistemáticos (es decir, el valor esperado de sus errores de previsión era nulos) ya que utilizaban toda la información disponible para predecir el futuro. Además, los economistas de las teorías RBG daban una enorme importancia a la calibración, simulación, estimación y evaluación cuantitativa de sus modelos.
} No obstante, presentaba grandes limitaciones, no solo de especificación (como las observadas antes) sino también cuando se intentaba contrastar los resultados con la realidad. En particular:\footnote{%
    Se ha tratado aquí un modelo RBC muy básico, pero existe una amplísima literatura de extensiones reales o walrasianas de este que busca mejorar el ajuste a la realidad observada del ciclo. Como extensiones walrasianas se puede citar la inclusión del trabajo indivisible de HANSEN (1985), la multiplicidad de sectores con shocks específicos (LONG y PLOSSER, 1983; LILIEN, 1982), la introducción de impuestos distorsionantes (BAXTER y KING, 1993; CAMPBELL, 1994), producción doméstica (BENHABIB, ROGERSON y WRIGHT, 1991; GREENWOOD y HERCOWITZ, 1991) o riesgo de ingresos idiosincrático no asegurable (KRUSELL y SMITH, 1998).

Las dos primeras limitaciones que se expondrán fueron contrastadas y demostradas por: COGLEY y NASON (1995) y ROTEMBERG y WOODFORD (1996)
}
\begin{lnum}
    \item \textbf{No persistencia} (se introduce ad-hoc: carece de microfundamentación). El modelo carece de un mecanismo válido de persistencia: la dinámica de la producción sigue muy de cerca a la de los shocks, cuando aquella debería durar más tiempo. De este modo, el modelo debe incorporar como un supuesto la persistencia del propio shock (que sigue procesos autorregresivos) para conseguir persistencia en la producción y el resto de las variables, que logra en detrimento de la microfundamentación del modelo;
    \item \textbf{Mínima volatilidad}. Las variaciones predecibles de la producción en el modelo RBC básico son muy inferiores a las encontradas en la realidad y con características muy distintas; y,
    \item \textbf{Neutralidad monetaria}. El modelo implica neutralidad monetaria como consecuencia de la inexistencia de fricciones/rigideces nominales. Sin embargo, el trabajo empírico realizado en tos 80 y 90 a través de la estimación de modelos SVAR indicaban que al menos en el corto plazo los instrumentos de política monetaria tenían impactos sobre la actividad real. Hoy en día existe una sólida evidencia de que los shocks monetarios tienen importantes efectos reales, luego se prescinde de una fuente identificada de fluctuaciones. Las explicaciones candidatas para ello descansan en el ajuste nominal incompleto de precios o salarios.
\end{lnum}

Así, el conflicto entre las predicciones teóricas y la evidencia empírica fueron la principal motivación del surgimiento de nuevas lineas de investigación que pudieran tener mejor validación empírica, en especial por el efecto real de la Política Monetaria.

La más célebre de estas lineas fue precisamente la NEK-2. 

\textbf{¿Qué queremos decir a lo largo de este apartado?} - Por eso, en este último punto presentaremos el que (con cierta generalidad) podríamos llamar \textbf{el modelo canónico de la NEK-2}, repasando sus caracteristicas esenciales. Veremos también que, aunque existen importantes diferencias en cuanto a las implicaciones de Política Económica de los modelos RBC y los NEK-2, ambos modelos comparten importantes características en cuanto a su microfundamentación: optimización y agentes racionales (HER), shocks estocásticos y restricciones intertemporales de recursos (empresas y hogares).\footnote{%
    Ambos tipos de modelos contienen hogares representativos de vida infinita que maximizan la suma esperada descontada de utilidades (del consumo y del ocio y eventualmente, de los saldos reales) sujeto a una restricción presupuestaria intertemporal y un gran número de empresas que con un acceso a una tecnología idéntica sujeta a shocks aleatorios. Además, como en la teoría RBC, en los modelos neo-keynesianos el equilibrio toma la forma de un proceso estocástico para todas las variables endógenas del modelo consistente con la toma de decisiones óptimas y bajo expectativas racionales de hogares y empresas.
} No obstante, incorporan también rígideces nominales y mercados imprfectos (competencia monopolística), supuestos que permiten que la transmisión de la Política Monetaria y/o Fiscal sobre los precios de los bienes o sobre los salarios no sea inmediata y tenga consecuencias reales.

\subsection{Modelo canónico: Derivación analítica}
Aunque no existe una único formulación del modelo de la NEK, puesto que existen numerosas extensiones que se usan frecuentemente, todas comparten la idea de que su resolución conduzca a un estilizado sistema de dos ecuaciones que determinan el comportamiento de todo el modelo:

\eqblock{\begin{aligned}
    &\text{Curva IS}:&&\boxed{x_t = \mathbb E_t[x_{t+1}] - \frac{1}{\sigma}\big(i_t - \mathbb E_t[\pi_{t+1}] - r^n_t\big)}\\[2em]
    &\text{Curva de PHILLIPS}:&&\boxed{\pi_{t} = \beta \mathbb E_t[\pi_{t+1}] + \kappa x_t}
\end{aligned}
}{Modelo canónico de la NEK-2. En base a GALI y MONACELLI (2005)}

Veamos cada componentes separadamente.

\subsubsection{Curva IS: Hogares}
La curva IS dinámica es una versión linealizada de la \textbf{ecuación de EULER de los hogares} y relaciona negativamente producción y brecha de tipos efectivos que distorsionan la elección consumo-ahorro de los hogares ($r_t=i_t-\mathbb E_t[\pi_{t+1}]\gtrless r_t^n$, $r_t^n=\rho$):

\eqblock{\begin{aligned}
    &\text{Curva IS}:&&\boxed{x_t = \mathbb E_t[x_{t+1}] - \frac{1}{\sigma}\big(i_t - \mathbb E_t[\pi_{t+1}] - r^n_t\big)}\\[2em]
    &\text{Donde,}
    &&\left\{\begin{aligned}
        &\sigma&&\equiv \text{Inversa de la elasticidad intertemporal de sustitución}\\[1em]
        &(\underbrace{i_t-\mathbb E_t[\pi_{t+1}]}_{r_t\,\,\equiv\,\,\text{\scriptsize Tipo de interés real}}-r_t^n)&&\equiv\text{Brecha entre el Tipo de interés real y natural}
    \end{aligned}\right.
\end{aligned}
}{Curva IS: Componentes. Modelo canónico de la NEK-2}

Su interpretación económica parte del comportamiento optimizador de hogares en sus decisiones de consumo, pues ante incrementos del tipo de interés el consumidor puede optar por aumentar el ahorro y reducir el consumo.\footnote{%
    Esta interpretación económica difiere de la de la IS tradicional de la SNC, en que esta relación inversa se debía al efecto negativo del tipo de interés sobre la inversión y por ende sobre la demanda agregada. 

Recuérdese que en este modelo neokeynesiano no tiene cabida la inversión por haber prescindido de la variable capital.
} La FUI objetivo, separable y aditiva, incorpora tanto la desutilidad del trabajo como la utilidad del consumo, que en este caso será una cesta de bienes diferenciados.\footnote{%
    Además, puede incorporarse el dinero en su función de utilidad. Introducido como \textit{money-in-the-utility} function (MIU, à la Sidrauski, 1967), para obtener una demanda del mismo, ya que el papel de la autoridad monetaria es clave en estos modelos.
} Con esto conseguimos incorporar la competencia monopolística originada en el lado de la demanda (lo que concuerda con la Teoría dominante y la evidencia empírica) al modelo pues al tratarse de bienes diferenciados (à la DIXIT-STIGLITZ), i.e. aunque existan muchas empresas éstas tendrán poder de mercado para fijar precios.

Por otro lado, su restricción presupuestaria está formada por: (i) un componente de riqueza (ya sea en saldos monetarios reales o en bonos); (ii) otro de renta; (iii) participación en los beneficios empresariales; y, (iv) en algunos casos, transferencias e impuestos (positivas o negativas, neutras o distorsionantes) del Gobierno.

En definitiva, se logra relacionar la brecha de producción (calculada en base a la producción potencial de la economía dados sus fundamentos económicos) con las distorsiones que inciden en las decisiones intertemporales de los consumidores, además de un componente forward looking ($\mathbb E_t[x_{t+1}]$) que otorga persistencia y suavización de las perturbaciones. 

\subsubsection{Curva de PHILLIPS: Empresas}
Por otro lado, la curva de PHILIPS de la NEK-2 (NKCP) procede de las \textbf{decisiones de precios óptimas de las empresas}, que relaciona positivamente la inflación presente con las expectativas de inflación futura y la brecha de producción o output gap: 

\eqblock{\begin{aligned}
    &\text{Curva de PHILLIPS}:&&\boxed{\pi_{t} = \beta \mathbb E_t[\pi_{t+1}] + \kappa x_t}\\[2em]
    &\text{Donde,}
    &&\left.\begin{aligned}
        &\kappa=\lambda(\sigma+\varphi)&&
        \left\{\begin{aligned}
            &\varphi\equiv \text{Inversa elasticidad oferta de empleo}\\[0.5em]
            &\underset{\text{\scriptsize $\Longrightarrow$ Problema del productor dinámico + Efectos PM}}{\lambda\equiv \text{Coeficiente de rigideces nominales}}\\[0.5em]
        \end{aligned}\right.\\[1em]
    \end{aligned}\right.\\[1em]
    &\left[\underset{\text{\scriptsize diferente en un contexto de economía abierta}}{\text{\scriptsize$\text{recordar,}\quad\overset{\text{\tiny proporcionales}}{mc_t\propto x_t}:mc_t=(\sigma+\varphi)x_t$}}\right]
\end{aligned}
}{Curva NKPC: Componentes. Modelo canónico de la NEK-2}

Dado que ya hemos supuesto que las empresas pueden tener poder de mercado ($mc_{EE}=-\mu$), al añadir rigidez en los precios ($\lambda$), tenemos todo lo que necesitamos para que la transmisión de la Política Monetaria tenga efectos reales:\footnote{%
    Recordar que la forma de introducir las rigideces tiene dos formulaciones alternativas popularmente utilizadas: (i) la que en esta versión se utiliza, donde las empresas están sujetas a restricciones sobre la frecuencia con la cual ajustan precios ante cambios en la demanda (ajuste escalonado de precios, entre los que destaca la fijación de precios à la Calvo (1983)32; y/o, (ii) las empresas están sujetas a costes de ajuste en los precios ante cambios en la demanda (costes de ajuste à la Rotemberg (1982), la versión dinámica de los costes de menú).

La misma clase de poder de mercado y fricciones pueden ser introducidas en el mercado de trabajo para obtener salarios rígidos.
} las empresas, en previsión de mayor inflación futura, decidan anticipar hoy estos mayores precios, debido a las dificultades para ajustar precios en el futuro cuando esa inflación esperada se haga efectiva. 

Además, nótese que por la introducción de precios escalonados, el problema de maximización de la empresa se hace verdaderamente dinámico, a diferencia del modelo de la RBC.\footnote{%
    Los beneficios en un período dependen de los precios fijados en periodos anteriores. Por ejemplo, en el caso de rigidez á la Calvo, en que sólo un porcentaje dado (y aleatorio) de las empresas puede ajustar precios en cada periodo, cuando a la empresa le toque ajustar precios, lo hará no solo pensando en el precio óptimo para ese período sino también descontando que en periodos futuros es posible que no pueda ajustar.
} Por otro lado, como se prescinde del capital, la FPA suele incluir como único factor productivo el trabajo, que se puede también incorporar con imperfecciones de mercado. 

Y, por último, el componente forward looking ($\mathbb E_t[\pi_{t+1}]$) otorga la persistencia necesaria sin recurrir a prámetros no microfundamentados.

\subsubsection{Regla de Taylor: Autoridad Monetaria}
Todo lo anterior nos conduce a un resultado esencial: \textbf{la no-neutralidad de la Política Monetaria en el corto plazo}.\footnote{%
    Por ejemlpo, la presencia de rigideces nominales ($\lambda$) hace que cambios en el tipo de interés a corto plazo ($i_t$, producido por cambios en el tipo de interés al que la autoridad monetaria presta a la banca comercial) no sean ajustados por cambios en la inflación esperada conduciendo a variaciones en los tipos de interés reales ($r_t=i_t-\mathbb E_t[\pi_{t+1}]>r_t^n$). 

Esta brecha entre en los tipos de interés producirá cambios en el consumo y en la inversión y, como resultado, en el producto ($x_t$) y en el empleo ($\varphi$) ya que las empresas encuentran óptimo ajustar la cantidad de bienes ofertados al nuevo nivel de demanda. A pesar de que, en el largo plazo, todos los precios y salarios se ajustarán y la economía volverá a su equilibrio natural.
}

Por ello, es común describir una función de reacción al estilo de Regla de Tipo de interés (de Taylor, por ejemplo), que describa el comportamiento de la Autoridad Monetaria, pues su tarea es precisamente usar el tipo de interés nominal ($i_t$) para combatir la inflación y lograr cerrar dicho gap. Una forma general de esta función de reacción/función objetivo es:

\eqblock{\begin{aligned}
    &\text{Regla de TAYLOR}:&&\boxed{r_t=\phi_\pi\mathbb E[\pi_{t+1}]+\phi_{x}\mathbb E_t[x_{t+1}]}\\[2em]
    &\text{Donde,}
    &&\left.\begin{aligned}
        &(\phi_\pi,\phi_x)\in[0,1]:(1-\phi_\pi)=\phi_x&&\equiv \text{Parámetros de ponderación}
    \end{aligned}\right.\\[1em]
    &\text{¿Óptima?}\Longrightarrow &&\underset{\text{\scriptsize $\Longrightarrow$ elimina efectos rigideces nominales: no hay dispersón de precios}}{\boxed{\text{Obj.}\quad\pi_t=0}\qquad \text{p.q.}:\text{si }\rho=r_t^n,\,\,\boxed{r_t=\rho}}
\end{aligned}
}{Regla de Taylor: Función de reacción de la Autoridad Monetaria. Modelo canónico de la NEK-2}

Sin entrar a analizar los diferentes objetivos que puede establecer una Autoridad Monetaria, podemos decir que, con cierto grado de consenso y como consecuencia de la especificación general del modelo la literatura considera que la regla óptima es el \textit{inflation targeting}, tal que $\mathbb E_t[\pi_{t+1}]=0$.\footnote{%
    De manera que el tipo de interés real efectivo sea igual al tipo de interés natural de la economía, el factor subjetivo de descuento.

Esto es así porque el coste marginal real de la economía es precisamente $\widehat{mc}_t=mc_t-mc$, donde $mc$ es el coste marginal real de la empresa en Estado Estacionario, que está directamente relacionado con el poder de mercado que este ostenta, y en consecuencia el mark-up: $mc=-\mu=-\log{\frac{\varepsilon}{\varepsilon-1}}$.

La idea central es que, en este entorno, la inflación introduce una dispersión ineficiente de precios relativos entre empresas que producen bienes diferenciados, y esa dispersión actúa como una cuña real que aleja la economía de su asignación eficiente. En cada periodo solo una fracción de empresas puede reoptimizar su precio, mientras que el resto mantiene precios fijados en el pasado. Si hay inflación positiva o negativa, los precios relativos de esas empresas empiezan a divergir: las que no pueden ajustar ven cómo su precio relativo se vuelve demasiado bajo o demasiado alto en términos reales. Esa dispersión implica que bienes muy similares se venden a precios distintos, lo cual genera una mala asignación de la producción entre firmas, porque la demanda se desvía hacia las que tienen precios relativamente más bajos aunque no sean las más eficientes. Esta ineficiencia no existiría en un mundo sin rigideces nominales.
} 

\paragraph{¿Divina Coincidencia? Política Monetaria y \textit{Shocks}}
De hecho, siendo este el caso y teniendo la Autoridad Monetaria el suficiente grado de credibilidad, llegamos a un conclusión aún más optimista. 

Debido a la interrelación entre inflación y optimización de las empresas, siendo la primera una distorsión sobre la segunda que se ve acentuada por la existencia de rigideces nominales; si la Autoridad Monetaria logra \textbf{convencer a los agentes de su compromiso inflacionista}, entonces logrará estabilizar a la vez la inflación y el output, eliminado las distorsiones reales que genera la presencia de rigideces nominales y alcanzando lo que se conoce como \textbf{Divina Coincidencia}. 

\textbf{INTRODUCIR GRÁFICO EN EL QUE SE ALCANZA LA "DIVINA COINCIDENCIA"}

No obstante, lo más común en el marco de la Teoría de Ciclos de la NEK-2 es, precisamente, generar distorsiones estocásticas que impidan alcanzar esta estabilización simultanea de manera que se introduzca mayor diversidad de shocks que puedan llegar a desencadenar el ciclo.\footnote{%
    Tras un laborioso (y a veces insospechado) proceso de sustituciones y combinaciones de ecuaciones y log-linealización de estas. Por razones de extensión no se incorpora en el tema un modelo más o menos formalizado, pero es altamente recomendable consultar la especificación formal completa de este o algún otro modelo completo y estar mínimamente familiarizado con ella. En Walsh (2017) se ofrece una versión completa relativamente simple y bien explicada de un modelo neokeynesiano. Igualmente resulta muy revelador inspeccionar algún modelo más sofisticado, como el Smets-Wouters, a fin de tener una idea informada del grado de complejidad que pueden alcanzar.

Nótese que, a pesar de incorporar fricciones que impidan y retrasen en el tiempo el ajuste nominal completo, se siguen imponiendo perturbaciones autorregresivas.
} Por lo que, en el contexto de la Teoría de Ciclos, en la mayoría de modelos se introducen perturbaciones estocásticas:

\eqblock{\begin{aligned}
    &\text{Curva IS}:&&\boxed{x_t = \mathbb E_t[x_{t+1}] - \frac{1}{\sigma}\big(i_t - \mathbb E_t[\pi_{t+1}] - r^n_t\big)+u_t^{IS}}\\[2em]
    &\text{Curva de PHILLIPS}:&&\boxed{\pi_{t} = \beta \mathbb E_t[\pi_{t+1}] + \kappa x_t+u_t^{\pi}}\\[2em]
    &\text{Regla de TAYLOR}:&&\boxed{i_t=r_t+\phi_\pi\mathbb E[\pi_{t+1}]+\phi_{x}\mathbb E_t[x_{t+1}]+u_t^{MP}}\\[1em]
    &\text{donde,}&&u_t^i\sim N(0,\sigma_{i})
\end{aligned}
}{Modelo canónico de la NEK-2: Perturbaciones estocásticas. En base a GALI y MONACELLI (2005)}

De esta forma, a diferencia del caso RBC, las perturbaciones o shocks tiene \textbf{múltiples orígenes} (aunque, a menudo, desconocidos). La \textbf{propagación} se debe, al igual que en la teoría RBC, al carácter de equilibrio general del modelo (la interdependencia entre las decisiones de empresas y consumidores). Y en cuanto a la \textbf{persistencia}, el protagonismo es de las (diversas) fuentes de \textbf{rigidez nominal}:
\begin{lnum}
    \item \textbf{Precios}. Limitar el ajuste nominal por un lado atenúa la incidencia del shock puntual inicial, pero por otro favorece a su persistencia, pues lleva a que su efecto se reparta en múltiples períodos posteriores.\footnote{
    Cabe destacar que el ajuste nominal incompleto en un contexto dinámico no es una renuncia al ajuste, sino en general (dependerá de la fuente concreta de rigidez) un ajuste pospuesto o distribuido en múltiples períodos.

    \textbf{OJO}: No obstante, también cabe recordar que las rigideces nominales no suelen ser suficientes para obtener un nivel de persistencia coherente con la observada, de ahí que las perturbaciones introducidas deban ser en ciertos casos, como en el modelo RBC, procesos autorregresivos. Sin embargo, muchas extensiones del modelo canónico introducen sofisticaciones que otorguen al modelo mecanismos internos de persistencia, evitando así imponerla ad hoc con perturbaciones autorregresivas.
    } ; y,
    \item \textbf{Política Monetaria}. Y en segundo lugar, la presencia del sector público puede ser relevante pues, además de contener una fuente de shocks nominales, sus políticas estabilizadoras buscarían recuperar el equilibrio sin rigideces, con lo que con su actuación contrarrestarían en parte dichas rigideces nominales (mediante sobrerreacciones y generando un trade off acentúado entre $x_t$ y $\pi_t$). 
\end{lnum}

\subsection{Implicaciones de Política Económica}
En definitiva, todas las modificaciones llevadas a cabo por la NEK-2 justifican la intervención del policy maker ante una perturbación que desvíe a la economía de su equilibrio: la respuesta de la economía a los diferentes shocks o ciclo será \textbf{ineficiente} y la no neutralidad de la Política Monetaria resultante de la presencia de rigideces nominales abre la puerta a que las intervenciones de la Autoridad Monetaria incrementen el bienestar en tanto reducirán las distorsiones existentes. Este adquiere un papel \textbf{protagonista en la suavización del Ciclo Económico}.

También cabe destacar que queda patente la influencia de la revolución metodológica iniciada por LUCAS, hasta el punto que podríamos considerar el modelo de la NEK-2 como un perfeccionamiento del modelo RBC. Dicho de manera muy simplificada, así como el modelo RBC es grosso modo un modelo RAMSEY con modificaciones, puede decirse que el modelo NEK-2 es un modelo RBC con modificaciones. Según el modelo de la NEK-2 el origen del ciclo sigue siendo exógeno y aleatorio, aunque ahora puede ser \textbf{real y/o nominal}. 

\subsubsection{Limitaciones: \textit{Forward guidance puzzle}}
No obstante, aunque este modelo presenta muchas virtudes, también presenta serios defectos. El más evidente, y que genera verdaderos dolores de cabeza dado el contexto actual es la conocida paradoja de la orientación futura o \textit{forward guidance puzzle}.

Dada la especificación del modelo, una vez que la Autoridad Monetaria adquiere la suficiente credibilidad como para influir en las expectativas de los agentes: el simple anuncio de un cambio de rumbo futuro de la Política Monetaria no solo tiene un \textbf{efecto inmediato sobre la producción}, independientemente de lo lejana que sea su implementación, sino que sorprendentemente su efecto acumulado sobre la producción es \textbf{proporcional a cuán anticipado es el anuncio}. Adicionalmente tiene un \textbf{efecto sobre la inflación desmesurado}, incluso mayor cuanto más lejana sea dicha implementación. Esta paradoja reposa en las debilidades más importantes del modelo: la optimización intertemporal plena en la demanda y la naturaleza absolutamente forward-looking de la demanda y de la fijación de precios. 

Además, desde una punto de vista empírico, el modelo carece de mecanismos que generen la suficiente persistencia de los shocks, su curva de PHILIPS implica que los costes de una desinflación son muy inferiores a los observados y su IS dinámica reposa sobre un comportamiento plenamente optimizador que parece no estar apoyado por la evidencia empírica. 

\clearpage
\section{Extensiones: Otras lineas de investigación}
\puntosclave{
    La crisis financiera global forzó un análisis detallado del sector financiero de una manera más realista, reconociendo sus imperfecciones. Si bien se ha ido incorporando a los modelos de ciclo económico, simultáneamente ha aparecido una literatura dedicada al ciclo financiero como ciclo paralelo, aunque no perfectamente sincronizado con el ciclo económico o de negocios visto hasta ahora. 
    
    El análisis del ciclo financiero presenta hasta la fecha un carácter eminentemente empírico, tratando de identificar los principales hechos estilizados que describen este ciclo financiero. 
    
    En efecto, el ciclo financiero parece poder describirse acudiendo a pocas variables (crédito y precios inmobiliarios), exhibe menor frecuencia (mayor amplitud) que el ciclo económico, está estrechamente ligado a la irrupción de crisis financieras, presentaría mayor predictibilidad (gracias a indicadores adelantados más fiables) y sería muy dependiente del régimen de políticas económicas (financieras, monetarias y reales) imperante. 
    
    El ciclo financiero presenta retos, tanto teóricos (se carece aún de un modelo de referencia que lo explique) como de políticas económicas (las políticas monetaria, fiscal y macroprudencial deben incorporar consideraciones a medio plazo, dada la mayor duración del ciclo financiero).
}

Como valoración de estos modelos de ciclo neokeynesianos, su evaluación puede situarse en algún punto entre dos opiniones extremas o polares:
\begin{la}
    \item \textbf{Optimista: Fundamentos ($\checkmark$)}. En un extremo, la visión optimista considera que estos modelos están suficientemente bien fundamentados en supuestos microeconómicos como para que los parámetros de que dependen tengan carácter estructural y, por tanto, son suficientemente realistas como para arrojar recomendaciones de Políticas en términos de bienestar;\footnote{%
        Se apunta que los modelos se construyen sobre fundamentos microeconómicos; que las simulaciones estimadas con los modelos captan razonablemente bien algunas características importantes de los ciclos; que muchas autoridades valoran los modelos lo suficiente como para apoyarse en sus predicciones y recomendaciones; que existe evidencia microeconómica para muchos de sus supuestos, y que su perfeccionamiento está avanzando con rapidez (ROMER, 2018).

    Se han desarrollado modelos DSGE más sofisticados que, aunque a costa de perder en transparencia y carecer de soluciones analíticas, son ampliamente utilizados por académicos, bancos centrales y gobiernos. Entre esas extensiones se pueden citar:
    \begin{la}
        \item \textbf{Oferta agregada}. La curva de PHILIPS neokeynesiana se puede extender en dos sentidos: (i) introduciendo inercia en la inflación; (ii) permitir también ajuste incompleto en salarios, típicamente à la Calvo, dando lugar a una curva de Philips neokeynesiana para la inflación salarial;
        \item \textbf{Demanda agregada}. La IS neokeynesiana presenta dos importantes limitaciones: (i) deja fuera la inversión, el gasto público y las exportaciones netas (las dos primeras son más habituales, la tercera se encuentra en modelos de la NOEM); y, (ii) implica grandes y rápidas respuestas de la producción ante shocks (incluso incorporando los otros componentes de la demanda agregada). Para ajustarse mejor a lo observado se suelen añadir ingredientes que retardan el ajuste, como hábitos en el consumo o costes de ajuste en la inversión; y,
        \item \textbf{Hipersensitivuty}. Para conseguir el exceso de sensibilidad del consumo a la renta observado en la realidad, se puede introducir que una fracción del consumo se determina por reglas fijas (sin optimizar, rules of thumb) o que un porcentaje de los hogares sufre restricciones de liquidez.
    \end{la}
    } y,
    \item \textbf{Pesimimista: Fundamentos ($\mathrm X$)}. En el otro extremo, la visión pesimista reconoce que estos modelos ofrecen una mala predicción en tiempo real (EDGE y GÜRKAYNAK, 2010) y, para ajustarse a los datos macroeconómicos, se les incorporan supuestos ad hoc (indexación, hábitos, costes de ajuste en la inversión) lo suficientemente alejados de la realidad como para que sus recomendaciones de política no sean totalmente fiables y sus predicciones probablemente fallen si el entorno macroeconómico cambia.
\end{la}


Como ya hemos repasado las virtudes de la Teoría de Ciclos de la NEK-2, veamos brevemente alguno de los argunentos que esgrimen sus oponentes. En este sentido se apunta a dos hecho: (i) hasta la reciente crisis el tratamiento de las imperfecciones del mercado de crédito era mínima; y, (ii)  la inclusión de muchas de las características del modelo no se justifica por evidencia microeconómica, sino por un deseo de explicar los hechos macroeconómicos.\footnote{%
    Como apunta ROMER (2018), en el análisis de los fundamentos microeconómicos del ajuste nominal incompleto en contexto dinámico, se aprecia una evidente falta de consenso acerca de la mejor manera de construir un modelo dinámico de fluctuaciones realista: el autor examina 7 alternativas distintas de ajuste dinámico de precios en su libro.
}

\textbf{¿Qué queremos decir?} Veremos brevemente las alternativas que existen para sortear la primera limitación, los posibles ciclos del mercado financiero, y algunos hechos estilizados que deberían poderse contrastar cuando se pretenda especificar un modelo macroeconómico de Ciclos.

\subsection{Ciclos Financieros: Imperfecciones del mercado de crédito}
Hasta la crisis de 2008-2009, los modelos DSGE en el marco de la NEK-2 no trataban el sector financiero, considerándolo como un velo que podía ser ignorado para comprender los ciclos económicos, o que a lo sumo contribuía a explicar parte de su persistencia (BORIO, 2012).

No obstante, considerando el papel determinante que jugó el sistema financiero en la misma resulta evidente que la necesaria introducción de las imperfecciones en el mercado del crédito. Es por ello que el interés de los economistas se dirigió a estudiar el sector financiero en el marco del análisis macroeconómico, pudiéndose observar dos grandes enfoques: (i) en el plano teórico, cada vez más modelos DSGE que incorporan el sector financiero de manera más sofisticada (DSGE aumentados con fricciones financieras), llegando incluso a explicar la aparición de crisis financieras endógenas (COLLARD, BOISSAY, GALÍ y MANEA, 2023);\footnote{%
    Para una visión de conjunto de estos modelos de crisis financieras puede consultarse el capítulo 10 de Romer (2018).
} y, (ii) estudio de los datos y características empíricas del sector financiero así como las series de variables financieras. 

\subsubsection{Acelerador Financiero: BERNANKE y GETLEY}
Una de las manera más célebres de incorporar los mercados financieros a la modelización es a través de la Teoría del Acelerador Financiero, propuesta por diferentes académicos, entre los cuales BERNANKE y GETLER.

Dicho de manera sencilla, el acelerador financiero es el \textbf{proceso mediante el cual un choque económico adverso puede verse amplificado por el empeoramiento de las condiciones del mercado financiero}. Este vínculo entre la economía real y el mercado financiero surge de los problemas de información asimétrica que existen entre prestatario y presetamista, cuando los primeros necesitan contar con financiación externa para cualquier empréstito, incluido el mantenimiento de su producción tras un shock adverso de la economía agregada.

De forma muy sencilla, podemos expresar las relaciones entre la economía real y el mercado de crédito como:

\eqblock{
\begin{aligned}
    &y_t^s(z_t):&&z_t=c_t+b_t\quad \text{donde,}
    \left\{\begin{aligned}
        &z_t\equiv \text{Insumos de producción}\\
        &c_t\equiv\text{Flujo de efectivo}\\
        &b_t\equiv\text{Endeudamiento}
    \end{aligned}\right.\\[1em]
    &\underset{\text{\scriptsize info. asim.: colateral}}{b_t\leq \mathbb E_t[v_{t+1}]}&&\Longrightarrow \boxed{z_t\leq c_t+\mathbb E_t[v_{t+1}]}\quad \text{donde,}\quad \mathbb E_t[v_{t+1}\,|\,y_t^s]
\end{aligned}
}{Acelerador financiero: Colateral. Teoría de los Ciclos financieros: BERNANKE y GETLER}

Por tanto, si se produce una caída de los precios de los activos ($v_t$), se deterioran los balances de las empresas y su patrimonio neto.
El consiguiente deterioro de su capacidad de endeudamiento ($b_t$) tiene un impacto negativo en su inversión y/o producción ($y_t^s(z_t)$). La disminución de la actividad económica reduce aún más los precios de los activos, lo que genera un ciclo de retroalimentación de caída de estos precios, deterioro de los balances, endurecimiento de las condiciones de financiación y disminución de la actividad económica.

\paragraph{Implicaciones de Política Económica: Política Monetaria}
Este círculo vicioso generado por el acelerador financiero consiste entonces en un bucle de retroalimentación financiera que, a partir de un pequeño cambio en los mercados financieros, es, en principio, capaz de producir un gran cambio en las condiciones económicas.

De esta forma, la autoridad monetaria debe intervenir relajando las condiciones crediticias y asegurándose que el deterioro de los activos tenga el menor impacto posible en la economía real. Curiosamente, tenemos dos ejemplo históricos polares tremendamente favorables hacia esta teoría:
\begin{la}
    \item \textbf{Gran depresión}. Muchos autores, entre los cuales destaca el estudio empírico realizada por FRIEDMAN y SCHWARTZ, destacan el papel de la Reserva Federal en la Gran Depresión generada por el crack bursátil de 1929. La estricta Política Monetaria que siguieron las autoridades acentúo las condiciones de crédito y acabo conduciendo a la economía a una depresión profunda de la que tardó mucho tiempo en salir. Es posible, argumentan algunos, que de haber reaccionado rápido y relajando los tipos, la economía real nunca hubiese experimentado tal nivel de recesión; y,
    \item \textbf{Crisis financiera (2007-2009)}. En contraste, la laxitud que caracterizó a la Política Monetaria inmediatamente después del estallido de la crisis financiera en 2007, permitió a la economía recuperar los niveles previos en un tiempo relativamente corto, entre 2-3 años.
\end{la}


\subsubsection{Burbujas financieras: BLANCHARD}
Otra linea de investigación complementaria, no alternativa, es la propuesta por BLANCHARD, gran exponente del marco de la NEK-2. Este autor trata de explicar la volatilidad en los mercados financieros como causa de la posible existencia de burbujas especulativas: una variable cuyo valor se desvía sistemáticamente y de manera progresiva (al alza o a la baja) de su valor fundamental. 

\eqblock{
\begin{aligned}
    \boxed{v_t=\sum_{s=0}^\infty \beta^s \underset{\text{\scriptsize valor fundamental esperado}}{\mathbb E_t[k_{t+s}]}+\varepsilon_t}\qquad \text{donde,}\;\;\underset{\text{\scriptsize valor burbuja generado por los mercados financieros}}{\varepsilon_t=\lim_{s\to\infty}\beta^s\mathbb E_t[v_{t+s}]}
\end{aligned}
}{Burbujas especulativas: Asuencia condición de transverslidad (BLANCHARD). Desarrollos recientes}

Por tanto, el valor actual del colateral (activo) será una función del valor esperado descontado al presente de la evolución futuro de sus valores fundamentales y un término de \textit{error} que en este caso viene determinado por el valor descontado al presente de la expectativas sobre la evolución del colateral futuro. 

\paragraph{Implicaciones de Política Económica: Supervisión}
Hemos hecho énfasis en que se trata de un marco completario y no alternativo porque la propuesta de BLANCHARD permite endogeneizar el fenómeno de \textbf{burbujas financieras}, situaciones en las que el valor del colateral no viene dado por lo que caracteriza sus fundamentos (inversores fundamentalistas que serían el primer término) sino que juegan un papel esencial los comportamientos especulativos que absorbe el segundo término (los chartistas), desde una perspectiva \textit{forward looking}. 

Por todo esto, las burbujas especulativas y su consecuente impacto sobre la economía real surgen debido a las imperfecciones de determinados mercados que acaban generalizándose al conjunto de la economía.

Esta línea de investigación justificaría por tanto la intervención del Policy Maker tanto desde una perspectivo directa, a través de la flexibilización de las condiciones crediticias, como desde una perspectiva indirecta, a través de la supervisión de los mercados financieros.

\subsection{Hechos Estilizados}
Hasta hora hemos expuesto el marco analítico a través del cual podemos estudiar el Ciclo o, más concretamente, sus consecuencias. Durante este, todas las variables estudiadas experimentan fluctuaciones de manera más o menos sincronizada, obviamente permitiendo que algunas variables se comporten de manera contracíclica (como el desempleo) o que tengan una naturaleza adelantada o retardada respecto al PIB. Pero en conjunto todas las variables comparten un mismo ciclo de fluctuación, es decir, una misma duración (o alternativamente frecuencia) del episodio entre pico y valle. 

\subsubsection{Correlación entre variables: Reales y Nominales}
No obstante, cabe hacer hincapíe en algunos hechos estilizados que se han venido repitiendo a lo largo de la historia económica:

\eqblock{
\begin{aligned}
    &\text{Real}:
    &&\left\{\begin{aligned}
        &Corr(x_t,c_t)>0,\quad Var(x_t)>Var(c_t)\\[1em]
        &Corr(x_t,i_t^o)>0,\quad Var(x_t)<Var(i_t^o)
        \quad \text{pero,}\left\{\begin{aligned}
            &Var(x_t)<Var(i_t^{D})\\[0.5em]
            &Var(x_t)\geq Var(i_t^{ND})
        \end{aligned}\right.\\[1em]
        &Corr(x_t,n_t)>0,\quad \underset{\text{\scriptsize dependerá de marco institucional}}{Var(x_t)\gtrless Var(n_t)}
    \end{aligned}\right.\\[1em]
    &\underset{\text{\scriptsize mcdos. financieros}}{\text{Nominal}:}
    &&\left\{\begin{aligned}
        &Corr(p_t^v,b_t)\quad\underset{\text{\scriptsize colateral y crédito, respectivamente}}{\text{Precios y volúmen}}\\[0.5em]
        &\underset{\text{\scriptsize menor frecuencia pero mayor impacto, colas gruesas}}{\rho_{b}(1)>\rho_x(1),\quad \rho_{p^v}(1)<\rho_{x}(1)\quad \text{pero},\,\,
        \left\{\begin{aligned}
            &Kurt(b_t)>Kurt(x_t)\\[0.5em]
            &Kurt(p_t^v)>Kurt(x_t)
        \end{aligned}\right.}\\[0.5em]
        &\text{Momento MINSKY ($\checkmark$)}: Corr(\sigma_b^2,\sigma_{p_v}^2)\,\text{y},\, Corr(b_t,\sigma_b^2)>0\\[0.5em]
        &Corr(\mathbb E_t[(b_t-\mu_b)^4],b_t)>0\quad \wedge\quad Cov(\mathbb E_t[(p_t^v-\mu_{p^v})^4],b_t)>0
    \end{aligned}\right.
\end{aligned}
}{}

No obstante, podría pensarse que, junto a este ciclo de negocios que afecta a un conjunto dado de variables, podría coexistir simultáneamente uno o incluso varios ciclos que afectaran a diferentes conjuntos de variables. Esto es lo que plantea el análisis empírico de los ciclos financieros, para cuya visión panorámica se seguirá el trabajo seminal de BORIO (2012), que reúne las principales ideas aprendidas del estudio del ciclo financiero. 

Como principales características estilizadas del ciclo financiero, Borio (2012) considera:
\begin{lnum}
    \item \textbf{Ciclo financiero: volumen y precio}. El ciclo financiero puede describirse de la manera más parsimoniosa en términos de dos variables, una de cantidad y una de precio: crédito (volúmenes) y precios. Estas dos variables parecen ofrecer el conjunto más pequeño de variables necesarias para replicar la interacción reforzada mutua entre restricciones financieras (crédito) y percepciones de valor y riesgo (precios), capturando los elementos esenciales del vínculo entre ciclo financiero, ciclo económico o de negocios y crisis financieras;\footnote{%
        Estas variables covarían de manera similar debido a la importancia del crédito en el sector de la construcción. Esto no implica que muchos análisis puedan basarse en conjuntos de series más reducidos (solo crédito) o más amplios (tipos de interés, volatilidad, primas de riesgo, tasas de impagos, morosidad, etc.).

    Además, es posible medir el aumento del riesgo de crisis financiera en tiempo real y con bastante exactitud (a diferencia del ciclo económico). Los indicadores adelantados de crisis financieras más prometedores se basan en desviaciones positivas simultáneas de la ratio de crédito privado sobre el PIB y de precios de los activos (especialmente inmobiliarios) respecto a sus medias históricas. El primero, una medida aproximada del endeudamiento, ofrece una visión de la capacidad de absorción de pérdidas del sistema (menor cuanto mayor sea el endeudamiento). Y el segundo, una medida aproximada de la probabilidad y magnitud de la ulterior reversión de los precios a su norma, lo que pone a prueba dicha capacidad de absorción de pérdidas.
    }
    \item \textbf{Ciclo financiero vs. ciclo real}. El ciclo financiero exhibe mucha menor frecuencia (es decir, mucha mayor duración y amplitud) que el ciclo de negocios tradicional;\footnote{%
        En efecto, el ciclo económico tiende unas frecuencias de entre 1 y 8 años, mientras la duración media del ciclo financiero en una muestra de siete países industrializados de la OCDE desde los años 60 es de alrededor de 16 años.
    }
    \item \textbf{Crisis financieras: Momento Minsky}. Los picos del ciclo financiero están estrechamente relacionados con el estallido de crisis financieras (crisis bancarias sistémicas).\footnote{%
        En la muestra antes citada, todas las crisis financieras de origen interno o nacional ocurrieron en o cerca del pico del ciclo financiero. Y esto explica una importante regularidad empírica: las recesiones (del ciclo económico) que coinciden con la contracción del ciclo financiero (haya o no crisis financiera) son especialmente severas (de media el PIB cae un 50\% más en estos casos).
    }
\end{lnum}

Y, por último, es evidente que el análisis económico carece de parámetros que representen el marco institucional de la economía. No obstante, la longitud y amplitud del ciclo financiero no son constantes, sino que dependen de las políticas económicas en vigor. Tres factores parecen ser importantes: el régimen financiero (la liberalización financiera debilita las restricciones de financiación, facilitando la plena interacción percepciones actitudes-condiciones financieras), el régimen monetario (si se restringe al control de la inflación a corto/medio plazo no actuará si el boom financiero se produce en entornos de inflación baja y estable) y el régimen de la economía real (factores de oferta como la globalización, que incrementan el crecimiento potencial y con ello el margen para expansiones del crédito y del precio de activos, a la vez que reducen la inflación). La evidencia empírica lo apoya: desde mediados de los 80, la longitud y amplitud del ciclo financiero se han incrementado a la vez que se liberalizaba el sector financiero, se instalaban exitosos regímenes monetarios de control de la inflación y se integraban más países (especialmente China) en el sistema global de comercio. En los siete países OCDE de la muestra, el ciclo financiero pasó de 11 años antes de 1998 a 20 años posteriormente.



Desde el punto de vista analítico, el ciclo financiero plantea retos para su explicación a través de modelos teóricos. En este sentido, tres serían los rasgos esenciales que deberían ser explicados por un modelo, según Borio (2012): a) la expansión financiera no debería solo preceder a la contracción, sino también ser su causa; b) un tratamiento explícito del sobreendeudamiento y de los excesivos stocks de capital como desequilibrios durante las contracciones, y c) una definición más rica de producción potencial, no tanto como producción no inflacionaria (la habitual) sino como producción sostenible (la inflación tradicionalmente ha sido usada como LA variable que refleja la información acerca de la diferencia entre el nivel de producción real y el potencial u output gap, pero en la práctica una inflación estable puede coexistir con una senda inestable debido a desequilibrios financieros y las distorsiones que estos ocultan en la economía real). Lograrlo requiere captar mejor los fallos de coordinación inherentes al ciclo financiero, y para ello se sugiere abandonar expectativas consistentes con el modelo (permitir que no haya consenso entre los agentes sobre cómo funciona la economía y cuál es el modelo correcto), incorporarpercepciones del riesgo u actitudes frente al riesgo que varíen sistemáticamente a lo largo del ciclo y, sobre todo, captar más profundamente la naturaleza monetaria de nuestras economías, en las que los intermediarios no asignan simplemente recursos reales sino que generan poder de compra (el crédito crea el depósito y no al revés) y en las que estos procesos interactúan con percepciones del valor débilmente ancladas, generando con ello inestabilidad. En suma, se requeriría superar el enfoque basado en el equilibrio de los modelos DSGE para redescubrir el análisis del desequilibrio.

Asimismo, el ciclo financiero presenta retos de política económica, de modo que las políticas monetaria, fiscal y macroprudencial deberían tener un enfoque al medio plazo, dada esa mayor duración del ciclo financiero. El principio básico sería generar reservas durante el auge financiero para aprovecharlas durante la contracción, estabilizado así el sistema. Y en caso de que no pueda evitarse una recesión severa, las políticas deberían fomentar y apoyar el ajuste de los balances, en lugar de retrasarlo involuntariamente. Surge aquí el denominado debate “lean vs clean”, entre quienes propugnan que las políticas deben evitar a toda costa la crisis financiera (lean against it) y quienes consideran que lo más importante es resolver adecuadamente (clean) las crisis una vez se produzcan. Finalmente, a la vez que existe un ciclo financiero doméstico o nacional en cada país, se ha identificado igualmente un ciclo financiero global (Miranda-Agrippino y Rey, 2021). Esta versión global del ciclo financiero se centra en cómo las condiciones financieras globales afectan a las economías individuales. No obstante, ambos ciclos financieros presentan relaciones interesantes (Aldasoro, Avdjiev, Borio, Disyatat, 2020). Si bien ambos son “variaciones sobre un mismo tema” y comparten su naturaleza procíclica (amplificadora de fluctuaciones reales), presentan interesantes diferencias, como las variables relevantes (precios de activos financieros y flujos de capitales transnacionales para el ciclo global), regularidades empíricas (el ciclo global tiene menor duración, típicamente vinculado al ciclo de negocios tradicional) y las implicaciones de política económica (“dilemma not trilemma” de Rey (2015), señalando que, para un país concreto, una política monetaria independiente en el sistema financiero global solo puede alcanzarse restringiendo los flujos financieros internacionales reflejados en la cuenta financiera de la balanza de pagos, independientemente de un régimen cambiario fijo o flotante). No obstante, ambos ciclos financieros suelen converger en episodios de crisis.









\clearpage
\section{Conclusión} 

\subsection{Extensiones y relación con otras partes del Temario}


\subsection{Opinión}



\subsubsection{Idea final (Salida o cierra)}


\clearpage
\pagenumbering{gobble}
\MyBibliography
\MyBibItem{Autor}{Título}{Año}{DOI -- LINK}




%ANEXOS
\clearpage
\anexo{\textbf{Preguntas Test}}
%\input{\TemaCode/questions.tex} 


\clearpage
\anexo{Métodos de descomposición de variables}\label{anx:filtros_descomposicion}

Uno de los métodos más conocidos para separar ciclo de una tendencia no lineal es el filtro de HODRICK y PRESCOTT (1997). El principal inconveniente de este método (y otros similares) es su carácter estrictamente estadístico y ateórico.\footnote{%
    Por lo tanto, surge la cuestión de cómo valorar la bondad un determinado filtrado (en el caso del filtro Hodrick-Prescott, cuál es el valor “correcto” del parámetro clave lambda). Además, este método genera estimaciones imprecisas de la tendencia en los extremos de las series (por lo que conviene desechar los primeros y últimos valores) y no captura rupturas estructurales en las series (por ejemplo, debidas a reformas estructurales de calado).
} En realidad existen métodos alternativos más o menos sofisticados, como los modelos de descomposición estadística de BEVERIDGE-NELSON y de componentes no observables, o los filtros band-pass de BAXTERY KING (1999).


\end{document}